% !TeX program = xelatex
% Compile twice with XeLaTeX. All illustrations are native vector TikZ drawings.
\documentclass[12pt,a4paper,oneside]{report}
\XeTeXgenerateactualtext=1
\usepackage[margin=23mm,top=24mm,bottom=24mm,headheight=27pt]{geometry}
\usepackage{amsmath,amssymb,mathtools,bm}
\usepackage{booktabs,array,longtable,tabularx}
\usepackage[table]{xcolor}
\usepackage{tikz}
\usetikzlibrary{arrows.meta,calc,patterns}
\usepackage[most]{tcolorbox}
\usepackage{titlesec,needspace}
\usepackage{hyperref}
\usepackage{bookmark}
\usepackage{xepersian}
\settextfont[Scale=1.06,BoldFont=Amiri-Bold.ttf]{Amiri-Regular.ttf}
\setlatintextfont[BoldFont=texgyretermes-bold.otf,ItalicFont=texgyretermes-italic.otf,BoldItalicFont=texgyretermes-bolditalic.otf]{texgyretermes-regular.otf}
\definecolor{navy}{HTML}{163747}
\definecolor{teal}{HTML}{127C80}
\definecolor{gold}{HTML}{B18542}
\definecolor{paper}{HTML}{F3F7F7}
\definecolor{muted}{HTML}{5E7178}
\definecolor{alert}{HTML}{A45538}
\hypersetup{unicode=true,colorlinks=true,linkcolor=navy,urlcolor=teal,
 pdftitle={روش‌های آماری ـ جزوهٔ بازنویسی‌شده},pdfauthor={علیرضا جهانبخش و احسان شهبازی},pdfsubject={احتمال، برآورد نقطه‌ای، فاصلهٔ اطمینان و آزمون فرض}}
\setlength{\parindent}{0pt}
\setlength{\parskip}{5pt plus 1pt}
\linespread{1.12}
\raggedbottom
\setcounter{tocdepth}{1}
\setcounter{secnumdepth}{2}
\allowdisplaybreaks[2]
\emergencystretch=1em
% Page styles use explicit boxes to keep bidi direction groups balanced.
\makeatletter
\def\ps@lecture{%
  \def\@oddhead{\vbox{\hbox to\textwidth{\small
    \rl{روش‌های آماری}\hfil\rl{\leftmark}}%
    \kern4pt\hrule height .35pt}}%
  \let\@evenhead\@oddhead
  \def\@oddfoot{\hfil\thepage\hfil}%
  \let\@evenfoot\@oddfoot}
\makeatother
\pagestyle{lecture}
\renewcommand{\chaptermark}[1]{\markboth{#1}{}}
\titleformat{\chapter}[display]{\color{navy}\bfseries\Huge}{\large فصل \thechapter}{8pt}{}[\vspace{4pt}\color{gold}\titlerule]
\titlespacing*{\chapter}{0pt}{0pt}{18pt}
\titleformat{\section}{\Large\bfseries\color{navy}}{\thesection}{.6em}{}
\titleformat{\subsection}{\large\bfseries\color{teal}}{\thesubsection}{.6em}{}
\tcbset{enhanced,breakable,boxrule=.45pt,arc=1.5mm,
 left=3mm,right=3mm,top=2mm,bottom=2mm,before skip=9pt,after skip=9pt,
 fonttitle=\bfseries,coltitle=white,colback=paper,colframe=teal}
\newtcolorbox{definition}[1]{breakable=false,title={تعریف | #1},colframe=navy,colback=navy!3!white}
\newtcolorbox{result}[1]{breakable=false,title={نتیجهٔ کلیدی | #1},colframe=teal,colback=teal!4!white}
\newtcolorbox{important}[1]{breakable=false,title={دقت کنید | #1},colframe=gold,colback=gold!6!white}
\newcounter{ex}[chapter]
\newenvironment{example}[1]{\refstepcounter{ex}\begin{tcolorbox}[breakable=false,title={مثال \thechapter.\arabic{ex} | #1},colframe=muted,colback=white]}{\end{tcolorbox}}
\newcommand{\edit}[1]{\par\begingroup\small\color{alert}\textbf{تکمیل یا اصلاح علمی.} #1\par\endgroup}
\newcommand{\source}[1]{\par{\small\color{muted} مبنا: صفحه‌های #1 فایل دست‌نویس.}\par}
\newcommand{\E}{\mathbb E}
\newcommand{\Var}{\operatorname{Var}}
\newcommand{\Cov}{\operatorname{Cov}}
\newcommand{\MSE}{\operatorname{MSE}}
\newcommand{\Bias}{\operatorname{Bias}}
\newcommand{\SE}{\operatorname{SE}}
\newcommand{\iid}{\overset{\mathrm{iid}}{\sim}}
\newcommand{\N}{\mathcal N}
\newcommand{\R}{\mathbb R}
\newcommand{\Prob}{\mathbb P}
\newcommand{\sol}{\par\textbf{حل.} }
\newcommand{\figcap}[1]{\par\smallskip{\small\color{muted}#1}\par}
\begin{document}
\begin{titlepage}
\begin{tikzpicture}[remember picture,overlay]
\fill[navy] (current page.north west) rectangle ([yshift=-90mm]current page.north east);
\fill[gold] ([yshift=-90mm]current page.north west) rectangle ([yshift=-92mm]current page.north east);
\draw[teal!22,line width=.5pt] ([xshift=23mm,yshift=32mm]current page.south west)--([xshift=-23mm,yshift=32mm]current page.south east);
\end{tikzpicture}
\vspace*{5mm}
\begin{center}
\textcolor{white}{{\large جزوهٔ درسی و آموزشی}}\\[12mm]
\textcolor{white}{{\fontsize{36}{46}\selectfont\bfseries روش‌های آماری}}\\[5mm]
\textcolor{white!85!navy}{{\Large از احتمال تا استنباط آماری}}
\end{center}
\vspace{24mm}
\begin{center}
\begin{LTR}
\begin{tikzpicture}[x=1.65cm,y=2.7cm]
\fill[teal!13] plot[domain=-1.96:1.96,samples=90] (\x,{exp(-\x*\x/2)}) -- (1.96,0)--(-1.96,0)--cycle;
\draw[navy,line width=1.2pt] plot[domain=-3.5:3.5,samples=130] (\x,{exp(-\x*\x/2)});
\draw[muted,-{Latex}] (-3.7,0)--(3.8,0);
\draw[gold,dashed] (-1.96,0)--(-1.96,.1465) (1.96,0)--(1.96,.1465);
\node[navy] at (0,.42) {$1-\alpha$};
\node[below,muted] at (-1.96,0) {$-z_{1-\alpha/2}$};
\node[below,muted] at (1.96,0) {$z_{1-\alpha/2}$};
\end{tikzpicture}
\end{LTR}
\vspace{9mm}
{\large احتمال و توزیع‌ها \quad | \quad برآورد نقطه‌ای}\\[3mm]
{\large فاصله‌های اطمینان \quad | \quad آزمون فرض آماری}\\[12mm]
{\color{teal}\large بازنویسی علمی و حروف‌چینی از ۴۵ صفحهٔ دست‌نویس}\\[3mm]
{\color{muted} همراه با اثبات‌ها، مثال‌های حل‌شده و نمودارهای بازطراحی‌شده}\\[9mm]
{\color{navy}\large\textbf{استاد درس: دکتر سکینه دهقان}}\\[4mm]
{\color{navy}\large گردآورندگان جزوه: علیرضا جهانبخش و احسان شهبازی}
\end{center}
\vfill
{\small\color{muted} این نسخه بر اساس فایل ارسالی تهیه شده است. بازسازی‌های مؤثر بر معنا و اصلاحات علمی در متن و پیوست مشخص شده‌اند.}
\end{titlepage}
\pagenumbering{roman}
\chapter*{راهنمای استفاده از جزوه}
\addcontentsline{toc}{chapter}{راهنمای استفاده از جزوه}
این جزوه، بازنویسی محتوای آموزشی تمام ۴۵ صفحهٔ فایل ارسالی است. ترتیب چهار فصل اصلی حفظ شده، اما مرزبندی بخش‌ها برای پیوستگی مطالعه تنظیم شده است. تعریف‌ها، استدلال‌ها، مثال‌ها و روابط تکرارشونده به زبانی یکدست نوشته شده‌اند؛ حل مثال‌هایی که در دست‌نویس تنها صورت آن‌ها آمده، تکمیل شده است.

همهٔ شکل‌ها از نو و به‌صورت برداری رسم شده‌اند. باکس‌های سرمه‌ای به تعریف‌ها، باکس‌های سبزآبی به نتایج اصلی، و باکس‌های طلایی به شرایط کاربرد و خطاهای متداول اختصاص دارند. یادداشت‌های «تکمیل یا اصلاح علمی» تغییرهای مؤثر نسبت به متن مبنا را توضیح می‌دهند.
\begin{important}{قراردادهای ثابت در سراسر جزوه}
پارامتر دوم $\N(\mu,\sigma^2)$ \textbf{واریانس} است. نماد $S^2$ همیشه واریانس نمونه با مخرج $n-1$ و $S$ انحراف معیار نمونه است. تمام چندک‌ها بر اساس \textbf{احتمال تجمعی از سمت چپ} تعریف می‌شوند:
\[
\Prob(Z\le z_p)=p,\quad
\Prob(T_\nu\le t_{p,\nu})=p,\quad
\Prob(Q_\nu\le\chi^2_{p,\nu})=p.
\]
در هر فرمول، فرض استقلال، نرمال بودن یا بزرگ بودن نمونه را پیش از استفاده بررسی کنید. علامت $\approx$ نشان‌دهندهٔ تقریب است و با برابری دقیق تفاوت دارد.
\end{important}
وقتی دادهٔ عددی یا عبارت دست‌نویس خوانش قطعی نداشته، خوانش انتخابی همراه با توضیح ذکر شده است. هیچ دادهٔ مبهمی بدون مشخص کردن فرض جایگزین نشده است. پیوست پایانی، نقشهٔ تطبیق با صفحات فایل مبنا و مهم‌ترین اصلاحات را در اختیار می‌گذارد.
\tableofcontents
\clearpage\pagenumbering{arabic}
\chapter{مقدمات احتمال و توزیع‌ها}
\source{۱ تا ۴}
\section{آزمایش تصادفی، فضای نمونه و متغیر تصادفی}
\begin{definition}{زبان احتمال}
\textbf{آزمایش تصادفی} فرایندی است که نتیجهٔ دقیق اجرای آن پیش از انجام معلوم نیست، هرچند مجموعهٔ نتایج ممکن را می‌شناسیم.

\textbf{فضای نمونه}، با نماد $\Omega$، مجموعهٔ تمام نتایج ممکن آزمایش است. هر زیرمجموعهٔ قابل اندازه‌گیری از آن یک \textbf{پیشامد} نام دارد.

\textbf{متغیر تصادفی} تابعی حقیقی و قابل اندازه‌گیری روی فضای نمونه است:
\[X:\Omega\longrightarrow\R.\]
\end{definition}
متغیر تصادفی گسسته، مقادیری در یک مجموعهٔ متناهی یا شمارا می‌گیرد. متغیر تصادفیِ دارای چگالی، مدلی برای کمیت‌هایی با مقادیر پیوسته است. در این جزوه، واژهٔ «پیوسته» در روابط چگالی به معنای \textbf{مطلقاً پیوسته} به کار می‌رود.
\edit{هر متغیر تصادفیِ بدون جرم نقطه‌ای الزاماً چگالی ندارد. روابط انتگرالی زیر مخصوص توزیع‌های دارای چگالی است.}
\section{تابع احتمال، چگالی و توزیع تجمعی}
\subsection{حالت گسسته}
تابع جرم احتمال $X$ به صورت
\[p_X(x)=\Prob(X=x)\]
تعریف می‌شود و باید شرایط زیر را داشته باشد:
\[p_X(x)\ge0,\qquad\sum_{x\in\mathcal S_X}p_X(x)=1.\]
در این رابطه $\mathcal S_X$ تکیه‌گاه $X$ است. تابع توزیع تجمعی، در همهٔ مدل‌ها، برابر است با
\[F_X(x)=\Prob(X\le x).\]
در حالت گسسته،
\[F_X(x)=\sum_{t\le x}p_X(t).\]
\subsection{حالت دارای چگالی}
\begin{result}{احتمال از مساحت زیر چگالی به دست می‌آید}
اگر $X$ چگالی $f_X$ داشته باشد،
\[
f_X(x)\ge0,\qquad\int_{-\infty}^{\infty}f_X(x)\,dx=1,
\qquad F_X(x)=\int_{-\infty}^{x}f_X(t)\,dt.
\]
برای $a<b$،
\[\Prob(a<X\le b)=F_X(b)-F_X(a)=\int_a^b f_X(x)\,dx.\]
همچنین $\Prob(X=b)=0$؛ ازاین‌رو باز یا بسته بودن دو سر بازه در محاسبهٔ احتمال تفاوتی ایجاد نمی‌کند.
\end{result}
هرجا مشتق وجود داشته باشد، $f_X(x)=F'_X(x)$؛ این برابری تقریباً همه‌جا برقرار است. تابع توزیع تجمعی نزولی نیست، راست‌پیوسته است و حدود آن در $-\infty$ و $+\infty$ به‌ترتیب صفر و یک‌اند.
\begin{center}
\begin{LTR}
\begin{tikzpicture}[x=1.45cm,y=1.8cm]
\fill[teal!20] plot[domain=-.7:1.5,samples=70] (\x,{exp(-\x*\x/2)})--(1.5,0)--(-.7,0)--cycle;
\draw[navy,thick] plot[domain=-3.2:3.2,samples=100] (\x,{exp(-\x*\x/2)});
\draw[-{Latex},muted] (-3.3,0)--(3.4,0) node[right] {$x$};
\draw[gold,dashed] (-.7,0)--(-.7,{exp(-.49/2)}) (1.5,0)--(1.5,{exp(-2.25/2)});
\node[below] at (-.7,0) {$a$};\node[below] at (1.5,0) {$b$};
\node[teal] at (.3,.35) {$\Prob(a<X\le b)$};
\end{tikzpicture}
\end{LTR}
\figcap{تعبیر هندسی احتمال در توزیع دارای چگالی؛ ارتفاع منحنی به‌تنهایی احتمال نیست.}
\end{center}
\section{توزیع مشترک و استقلال}
برای دو متغیر گسسته،
\[
p_{X,Y}(x,y)=\Prob(X=x,Y=y),\qquad
p_{X,Y}(x,y)\ge0,\qquad\sum_x\sum_y p_{X,Y}(x,y)=1.
\]
در حالت دارای چگالی مشترک،
\[
f_{X,Y}(x,y)\ge0,\qquad
\int_{-\infty}^{\infty}\!\int_{-\infty}^{\infty}f_{X,Y}(x,y)\,dx\,dy=1.
\]
توزیع‌های حاشیه‌ای با جمع یا انتگرال گرفتن نسبت به متغیر دیگر به دست می‌آیند؛ برای مثال
\[f_X(x)=\int_{-\infty}^{\infty}f_{X,Y}(x,y)\,dy.\]
\begin{definition}{استقلال دو متغیر تصادفی}
$X$ و $Y$ مستقل‌اند هرگاه برای تمام $x,y$،
\[F_{X,Y}(x,y)=F_X(x)F_Y(y).\]
در مدل گسسته، این شرط معادل $p_{X,Y}(x,y)=p_X(x)p_Y(y)$ است؛ در مدل دارای چگالی نیز $f_{X,Y}(x,y)=f_X(x)f_Y(y)$ تقریباً همه‌جا برقرار است.
\end{definition}
\section{امید ریاضی و خواص آن}
با فرض وجود امید ریاضی،
\[
\E(X)=\mu_X=
\begin{cases}
\displaystyle\sum_x x\,p_X(x),&\text{گسسته،}\\[3pt]
\displaystyle\int_{-\infty}^{\infty}x f_X(x)\,dx,&\text{دارای چگالی.}
\end{cases}
\]
برای تابع $g$، بدون نیاز به یافتن توزیع $g(X)$،
\[
\E[g(X)]=
\begin{cases}
\displaystyle\sum_x g(x)p_X(x),\\[3pt]
\displaystyle\int_{-\infty}^{\infty}g(x)f_X(x)\,dx.
\end{cases}
\]
\begin{result}{خطی بودن امید ریاضی}
برای ثابت‌های $a_i,b_i$ و متغیرهای با امید ریاضی متناهی،
\[
\E(a_iX_i+b_i)=a_i\E(X_i)+b_i,
\qquad
\E\left[\sum_{i=1}^n(a_iX_i+b_i)\right]
=\sum_{i=1}^n[a_i\E(X_i)+b_i].
\]
برای این خاصیت، استقلال متغیرها لازم نیست.
\end{result}
\section{واریانس، انحراف معیار و کوواریانس}
\begin{definition}{پراکندگی و تغییرات مشترک}
اگر گشتاور دوم متناهی باشد،
\[
\Var(X)=\sigma_X^2=\E[(X-\mu_X)^2]
=\E(X^2)-[\E(X)]^2,
\qquad \sigma_X=\sqrt{\Var(X)}.
\]
کوواریانس دو متغیر برابر است با
\begin{align*}
\Cov(X,Y)&=\E[(X-\mu_X)(Y-\mu_Y)]\\
&=\E(XY)-\E(X)\E(Y).
\end{align*}
\end{definition}
استقلال، با فرض وجود گشتاورهای لازم، به $\Cov(X,Y)=0$ منجر می‌شود؛ عکس این گزاره در حالت کلی درست نیست.

برای ثابت‌های $a,b,c,d$،
\[
\Var(X+b)=\Var(X),\qquad \Var(aX)=a^2\Var(X),
\qquad \Cov(aX+b,cY+d)=ac\Cov(X,Y).
\]
\begin{result}{واریانس ترکیب خطی}
\[
\Var\left(\sum_{i=1}^na_iX_i\right)
=\sum_{i=1}^na_i^2\Var(X_i)
+2\sum_{i<j}a_ia_j\Cov(X_i,X_j).
\]
اگر متغیرها دوبه‌دو ناهمبسته باشند، جملهٔ دوم صفر می‌شود. استقلال متقابل، شرط کافی برای این ساده‌سازی است.
\end{result}
\textbf{اثبات.} با قرار دادن $Y_i=X_i-\E(X_i)$ داریم
\begin{align*}
\Var\left(\sum_i a_iX_i\right)
&=\E\left[\left(\sum_i a_iY_i\right)^2\right]\\
&=\sum_i a_i^2\E(Y_i^2)+2\sum_{i<j}a_ia_j\E(Y_iY_j).
\end{align*}
این همان رابطهٔ بالا است. همچنین برای دو ترکیب خطی،
\[
\Cov\left(\sum_{i=1}^n a_iX_i,\sum_{j=1}^m b_jY_j\right)
=\sum_{i=1}^n\sum_{j=1}^m a_ib_j\Cov(X_i,Y_j).
\]
\subsection{نتایج مهم برای میانگین نمونه}
اگر $X_1,\ldots,X_n$ مستقل و هم‌توزیع با میانگین $\mu$ و واریانس $\sigma^2$ باشند،
\[
\bar X=\frac1n\sum_{i=1}^nX_i,\qquad
\E(\bar X)=\mu,\qquad \Var(\bar X)=\frac{\sigma^2}{n}.
\]
برای هر $j$،
\begin{align*}
\Cov(X_j,\bar X)
&=\frac1n\sum_{i=1}^n\Cov(X_j,X_i)
=\frac1n\Var(X_j)=\frac{\sigma^2}{n}.
\end{align*}
پس یک مشاهده و میانگین نمونه مستقل نیستند، زیرا خود مشاهده در ساخت میانگین به کار رفته است.
\section{توزیع‌های پرکاربرد}
در این بخش، تابع احتمال یا چگالی در خارج از دامنهٔ ذکرشده صفر است و $q=1-p$.
\subsection{توزیع‌های گسسته}
\textbf{یکنواخت گسسته.} اگر $X$ روی $\{x_1,\ldots,x_N\}$ احتمال‌های برابر داشته باشد،
\[
p_X(x_j)=\frac1N,\qquad
\E(X)=\frac1N\sum_{j=1}^N x_j,\qquad
\Var(X)=\frac1N\sum_{j=1}^N(x_j-\mu)^2.
\]
\textbf{برنولی.} $X\sim\operatorname{Bern}(p)$، با $0\le p\le1$:
\[
p_X(x)=p^xq^{1-x},\quad x=0,1;\qquad
\E(X)=p,\quad\Var(X)=pq.
\]
\textbf{دوجمله‌ای.} $X\sim\operatorname{Bin}(n,p)$:
\[
p_X(x)=\binom nx p^xq^{n-x},\quad x=0,\ldots,n;\qquad
\E(X)=np,\quad\Var(X)=npq.
\]
\textbf{پواسون.} $X\sim\operatorname{Pois}(\lambda)$، با $\lambda>0$:
\[
p_X(x)=\frac{\lambda^x e^{-\lambda}}{x!},\quad x=0,1,\ldots;
\qquad\E(X)=\Var(X)=\lambda.
\]
در دو مدل برنولی و دوجمله‌ای، نقاط مرزی $p=0,1$ مستقیماً از تعریف احتمال موفقیت تفسیر می‌شوند.
\subsection{توزیع‌های دارای چگالی}
\textbf{یکنواخت.} $X\sim U(a,b)$، با $a<b$:
\[
f_X(x)=\frac1{b-a},\quad a<x<b;\qquad
\E(X)=\frac{a+b}{2},\quad\Var(X)=\frac{(b-a)^2}{12}.
\]
\textbf{نمایی با نرخ $\lambda>0$.}
\[
f_X(x)=\lambda e^{-\lambda x},\quad x>0;\qquad
\E(X)=\frac1\lambda,\quad\Var(X)=\frac1{\lambda^2}.
\]
\textbf{گاما با شکل $\alpha>0$ و نرخ $\lambda>0$.}
\[
f_X(x)=\frac{\lambda^\alpha}{\Gamma(\alpha)}x^{\alpha-1}e^{-\lambda x},\quad x>0;
\qquad\E(X)=\frac\alpha\lambda,\quad\Var(X)=\frac\alpha{\lambda^2}.
\]
تابع گاما با $\Gamma(\alpha)=\int_0^\infty t^{\alpha-1}e^{-t}\,dt$ تعریف می‌شود و
\[
\Gamma(\alpha+1)=\alpha\Gamma(\alpha),\qquad
\chi^2_\nu\equiv\operatorname{Gamma}\left(\frac\nu2,\text{ نرخ }\frac12\right).
\]
\begin{important}{نرخ را با مقیاس اشتباه نگیرید}
اگر گاما با مقیاس $\beta=1/\lambda$ نوشته شود، امید و واریانس آن به‌ترتیب $\alpha\beta$ و $\alpha\beta^2$ است. تمام فرمول‌های این جزوه با پارامتر \textbf{نرخ} نوشته شده‌اند.
\end{important}
\chapter{برآورد نقطه‌ای و معیارهای کیفیت}
\source{۵ تا ۸}
\section{جامعه، نمونه، پارامتر و آماره}
\begin{definition}{مفاهیم پایهٔ استنباط}
\textbf{جامعهٔ آماری}: مجموعهٔ واحدهایی که می‌خواهیم دربارهٔ آن‌ها اطلاعات به دست آوریم.

\textbf{نمونه}: بخشی از واحدهای جامعه. نمونه‌گیری تصادفی به هر واحد، طبق طرح مشخص، امکان انتخاب می‌دهد؛ برابر بودن احتمال انتخاب، ویژگی نمونه‌گیری تصادفی ساده است.

\textbf{نمونهٔ تصادفی از یک توزیع}: متغیرهای $X_1,\ldots,X_n$ مستقل و هم‌توزیع‌اند؛ این فرض را با $\iid$ نشان می‌دهیم.

\textbf{پارامتر}: ویژگی عددی ثابت اما معمولاً نامعلومِ توزیع جامعه؛ مانند $\mu$ یا $\sigma^2$. مجموعهٔ مقادیر مجاز آن، فضای پارامتر $\Theta$ است.

\textbf{آماره}: تابعی از مشاهدات نمونه که به پارامتر نامعلوم وابسته نیست. آماره، پیش از مشاهدهٔ داده‌ها، متغیری تصادفی و دارای توزیع نمونه‌ای است.
\end{definition}
\textbf{برآوردگر} آماره‌ای مانند $T=T(X_1,\ldots,X_n)$ است که برای تخمین یک پارامتر به کار می‌رود. \textbf{برآورد} مقدار عددی $T(x_1,\ldots,x_n)$ پس از مشاهدهٔ داده‌هاست.
\begin{example}{توزیع نمونه‌ای میانگین}
جامعه‌ای شامل چهار مقدار $1,2,3,4$ داریم و با جایگذاری، دو مشاهده انتخاب می‌کنیم. هر زوج مرتب احتمال $1/16$ دارد. توزیع میانگین نمونه عبارت است از
\begin{center}\small
\begin{tabular}{c|ccccccc}
\toprule
$\bar x$&$1$&$1.5$&$2$&$2.5$&$3$&$3.5$&$4$\\
\midrule
$\Prob(\bar X=\bar x)$&$1/16$&$2/16$&$3/16$&$4/16$&$3/16$&$2/16$&$1/16$\\
\bottomrule
\end{tabular}\end{center}
برای نمونه، میانگین $2$ از زوج‌های $(1,3),(2,2),(3,1)$ حاصل می‌شود. میانگین جامعه $2.5$ و واریانس آن $1.25$ است؛ بنابراین $\E(\bar X)=2.5$ و $\Var(\bar X)=1.25/2=0.625$.
\end{example}
\section{نااریبی و اریبی}
\begin{definition}{برآوردگر نااریب}
برآوردگر $T$ برای $\theta$ نااریب است اگر برای تمام $\theta\in\Theta$،
\[\E_\theta(T)=\theta.\]
اریبی آن برابر است با
\[\Bias_\theta(T)=\E_\theta(T)-\theta.\]
\end{definition}
نااریبی به معنی برابری هر برآورد با پارامتر نیست؛ میانگین برآوردها در تکرار نمونه‌گیری برابر پارامتر است.
\begin{example}{نااریبی میانگین نمونه}
اگر $X_1,\ldots,X_n\iid\N(\theta,1)$، آنگاه
\[\E_\theta(\bar X)=\frac1n\sum_i\E_\theta(X_i)=\theta.\]
در واقع، این نتیجه برای هر نمونهٔ مستقل و هم‌توزیع با میانگین متناهی برقرار است و به نرمال بودن نیاز ندارد.
\end{example}
\begin{example}{چرا در واریانس نمونه بر $n-1$ تقسیم می‌کنیم؟}
جامعه میانگین $\mu$ و واریانس متناهی $\sigma^2$ دارد. تعریف کنید
\[\widehat\sigma_n^2=\frac1n\sum_{i=1}^n(X_i-\bar X)^2.\]
از همانی
\[\sum_i(X_i-\bar X)^2=\sum_i(X_i-\mu)^2-n(\bar X-\mu)^2\]
و $\E[(\bar X-\mu)^2]=\sigma^2/n$ نتیجه می‌شود
\[
\E\left[\sum_i(X_i-\bar X)^2\right]=(n-1)\sigma^2.
\]
پس
\[\E(\widehat\sigma_n^2)=\frac{n-1}{n}\sigma^2,\qquad
\Bias(\widehat\sigma_n^2)=-\frac{\sigma^2}{n}.\]
در مقابل، برای $n\ge2$،
\[\boxed{S^2=\frac1{n-1}\sum_i(X_i-\bar X)^2,\qquad \E(S^2)=\sigma^2.}\]
\end{example}
\section{خطای میانگین مربعات و کارایی نسبی}
\begin{result}{تجزیهٔ خطای میانگین مربعات}
برای برآورد $\theta$،
\[
\MSE_\theta(T)=\E_\theta[(T-\theta)^2]
=\Var_\theta(T)+[\Bias_\theta(T)]^2.
\]
در نتیجه، برای برآوردگر نااریب، $\MSE_\theta(T)=\Var_\theta(T)$ است.
\end{result}
\textbf{اثبات.} بنویسید $T-\theta=[T-\E(T)]+[\E(T)-\theta]$. پس از مجذور کردن و امید گرفتن، جملهٔ متقاطع به‌علت $\E[T-\E(T)]=0$ حذف می‌شود.

در این جزوه، $T_1$ را در مقدار پارامتر مورد نظر کاراتر از $T_2$ می‌نامیم اگر $\MSE_\theta(T_1)<\MSE_\theta(T_2)$. کارایی نسبی به صورت
\[
e(T_1,T_2)=\frac{\MSE_\theta(T_2)}{\MSE_\theta(T_1)}
\]
تعریف می‌شود؛ مقدار بیشتر از یک به سود $T_1$ است. برای دو برآوردگر نااریب، می‌توان نسبت واریانس‌ها را جایگزین کرد.
\begin{important}{مقایسه باید برای یک هدف مشترک باشد}
برای مقایسهٔ دو برآوردگرِ اریب، نسبت واریانس‌ها کافی نیست. هر دو خطا باید نسبت به \textbf{همان پارامتر هدف} محاسبه شوند. اگر ترتیب برتری با $\theta$ تغییر کند، هیچ‌یک در کل فضای پارامتر برتری یکنواخت ندارد.
\end{important}
\subsection{دو مثال از برآوردگرهای اریب}
\begin{example}{نمونهٔ دوجمله‌ای؛ خوانش حفظ‌شده از دست‌نویس}
فرض کنید $X_i\iid\operatorname{Bin}(k,\theta)$ و هدف برآورد $m=k\theta$ است. در متن مبنا دو آمارهٔ
\[A=\bar X,\qquad B=\frac{X_1-\bar X}{n}\]
برای مقایسه آمده‌اند. با $v=k\theta(1-\theta)$ داریم
\[
\E(A)=m,\quad\Var(A)=v/n,\qquad
\E(B)=0,\quad\Var(B)=\frac{v(n-1)}{n^3}.
\]
در محاسبهٔ واریانس $B$، کوواریانس $\Cov(X_1,\bar X)=v/n$ را لحاظ کرده‌ایم. بنابراین
\[
\MSE_m(A)=\frac vn,\qquad
\MSE_m(B)=\frac{v(n-1)}{n^3}+m^2.
\]
برای $0<\theta<1$،
\[
e(A,B)=\frac{n-1}{n^2}+\frac{nk\theta}{1-\theta}.
\]
\edit{صورت آمارهٔ دوم در دست‌نویس نامعمول است؛ خوانش «تفاضل تقسیم بر $n$» حفظ شده است. این آماره برای $k\theta$ نااریب نیست و مقایسهٔ صرفِ واریانس‌ها درست نخواهد بود.}
\end{example}
\begin{example}{نمونهٔ پواسون}
فرض کنید $X_i\iid\operatorname{Pois}(\theta)$ و $\theta>0$. دو برآوردگر $A=\bar X$ و $B=X_1+\bar X$ را برای $\theta$ مقایسه کنید.
\sol
\[
\E(A)=\theta,\quad\Var(A)=\theta/n,\qquad
\E(B)=2\theta,\quad\Bias_\theta(B)=\theta.
\]
به‌علت حضور $X_1$ در میانگین،
\[
\Var(B)=\theta+\frac\theta n+2\frac\theta n
=\theta\left(1+\frac3n\right).
\]
پس
\[
e(A,B)=\frac{\theta(1+3/n)+\theta^2}{\theta/n}
=n+3+n\theta>1.
\]
بنابراین $\bar X$ برای تمام $\theta>0$ کاراتر است.
\end{example}
\section{ترکیب خطی نااریب با کمترین واریانس}
\begin{example}{ضرایب نااریبی در یک نمونه}
اگر $X_i$ها میانگین $\mu$ داشته باشند، برای $T=\sum_i a_iX_i$ داریم
\[\E(T)=\mu\sum_i a_i.\]
پس برای نااریبی در تمام مقادیر $\mu$ لازم و کافی است $\sum_i a_i=1$. اگر علاوه بر آن، مشاهدات مستقل و واریانس‌های آن‌ها برابر $\sigma^2$ باشد،
\[\Var(T)=\sigma^2\sum_i a_i^2\ge\frac{\sigma^2}{n},\]
زیرا $\sum_i a_i^2\ge(\sum_i a_i)^2/n$. تساوی برای $a_i=1/n$ برقرار است؛ پس میانگین نمونه در این خانواده کمترین واریانس را دارد.
\end{example}
\begin{result}{وزن‌دهی معکوس واریانس}
اگر $T_1,T_2$ مستقل و نااریب برای $\theta$، با واریانس‌های مثبت $v_1,v_2$ باشند، ترکیب بهینهٔ نااریب آن‌ها عبارت است از
\[
T^*=\frac{v_2}{v_1+v_2}T_1+\frac{v_1}{v_1+v_2}T_2,
\qquad \Var(T^*)=\frac{v_1v_2}{v_1+v_2}.
\]
\end{result}
\textbf{اثبات.} شرط نااریبی $a_1+a_2=1$ است. تابع $a_1^2v_1+(1-a_1)^2v_2$ در $a_1=v_2/(v_1+v_2)$ کمینه می‌شود. اگر مانند یادداشت حاشیه‌ای متن مبنا $v_2=3v_1$ باشد، وزن‌ها به‌ترتیب $3/4$ و $1/4$ هستند.
\section{سازگاری}
\begin{definition}{نزدیک شدن احتمالی به پارامتر}
دنبالهٔ برآوردگرهای $T_n$ برای $\theta$ سازگار است اگر برای هر $\varepsilon>0$ و هر $\theta\in\Theta$،
\[
\Prob_\theta(|T_n-\theta|>\varepsilon)\longrightarrow0.
\]
این خاصیت را با $T_n\xrightarrow{P}\theta$ نشان می‌دهیم.
\end{definition}
با نامساوی مارکف روی $(T_n-\theta)^2$،
\[
\Prob_\theta(|T_n-\theta|>\varepsilon)
\le\frac{\MSE_\theta(T_n)}{\varepsilon^2}.
\]
پس $\MSE_\theta(T_n)\to0$ شرط کافی سازگاری است. به‌ویژه، اگر $T_n$ نااریب باشد و $\Var_\theta(T_n)\to0$، سازگار است. به‌طور کلی نیز کافی است اریبی و واریانس هر دو به صفر میل کنند.
\edit{میل کردن واریانس به صفر به‌تنهایی کافی نیست؛ مثلاً برآوردگر ثابت صفر برای پارامتر ناصفر، واریانس صفر دارد ولی سازگار نیست. شرط $\MSE\to0$ کافی است و ضرورت ندارد.}
\begin{example}{سازگاری میانگین}
اگر $X_i\iid\N(\theta,1)$، آنگاه $\MSE_\theta(\bar X_n)=1/n\to0$. بنابراین $\bar X_n\xrightarrow{P}\theta$.

همین استدلال برای $X_i\iid\operatorname{Pois}(\theta)$ با $\MSE=\theta/n$ و برای $X_i\iid\operatorname{Bin}(k,\theta)$، به‌منظور برآورد $\theta$ با $\bar X_n/k$، با $\MSE=\theta(1-\theta)/(nk)$ برقرار است.
\end{example}
\begin{important}{تمرین‌های ارجاع‌شده در متن مبنا}
در پایان این بخش، مثال‌های صفحه‌های ۱۰۵ تا ۱۱۲ و تمرین‌های صفحهٔ ۱۵۳، سؤال‌های ۱ تا ۵ کتاب معرفی شده‌اند. چون نام و متن کتاب در فایل موجود نیست، این ارجاع حفظ شده است و صورت سؤال‌های کتاب بازسازی نشده است.
\end{important}

\chapter{فاصله‌های اطمینان}
\source{۹ تا ۲۶}
\section{منطق فاصلهٔ اطمینان و کمیت محوری}
\begin{definition}{فاصلهٔ اطمینان}
فرض کنید داده‌ها از توزیعی با پارامتر ثابت و نامعلوم $\theta$ آمده‌اند. بازهٔ تصادفی $[L(X),U(X)]$ فاصلهٔ اطمینان با سطح $1-\alpha$ است اگر
\[\Prob_\theta\{L(X)\le\theta\le U(X)\}=1-\alpha.\]
در کاربردهای تقریبی، این احتمال تقریباً برابر $1-\alpha$ است. در برخی روش‌های محافظه‌کارانه نیز پوشش حداقل برابر این سطح است.
\end{definition}
پیش از نمونه‌گیری، دو کران تصادفی‌اند و پارامتر ثابت است. پس از مشاهدهٔ داده‌ها، بازهٔ عددی مشخص می‌شود. تعبیر تکرارپذیرِ اطمینان ۹۵ درصد این است که در تکرار نمونه‌گیری، تقریباً ۹۵ درصد بازه‌های ساخته‌شده پارامتر را می‌پوشانند.
\begin{important}{تعبیر احتمالی درست}
در استنباط فراوانی‌گرا، به پارامتر ثابت احتمال تصادفی نسبت نمی‌دهیم. عبارت «با اطمینان ۹۵ درصد، پارامتر در این بازه است» به پوشش روش ساخت بازه اشاره دارد؛ این عبارت، احتمال پسینِ پارامتر نیست.
\end{important}
\begin{definition}{کمیت محوری و چندک}
کمیت محوری تابعی مانند $Q=h(X_1,\ldots,X_n;\theta)$ است که توزیع آن به هیچ پارامتر نامعلوم وابسته نیست. خود $Q$ می‌تواند شامل پارامتر نامعلوم باشد و به همین دلیل الزاماً آماره نیست.

چندک $p$ برای توزیع پیوستهٔ $Q$، عدد $q_p$ است که $\Prob(Q\le q_p)=p$.
\end{definition}
برای ساخت فاصله، ابتدا
\[\Prob(q_{\alpha/2}\le Q\le q_{1-\alpha/2})=1-\alpha\]
را می‌نویسیم، سپس نامساوی را نسبت به $\theta$ حل می‌کنیم. جهت نامساوی‌ها هنگام معکوس‌کردن مقادیر مثبت اهمیت دارد.
\section{میانگین یک جامعهٔ نرمال؛ واریانس معلوم}
\source{۹ تا ۱۱}
اگر $X_1,\ldots,X_n\iid\N(\mu,\sigma^2)$ و $\sigma$ معلوم باشد، چون هر ترکیب خطی از متغیرهای \textbf{مشترکاً نرمال}، نرمال است،
\[
\bar X\sim\N\left(\mu,\frac{\sigma^2}{n}\right),\qquad
Z=\frac{\bar X-\mu}{\sigma/\sqrt n}\sim\N(0,1).
\]
استقلال متغیرهای نرمال در این مسئله، نرمال مشترک بودن را تضمین می‌کند. از تقارن $z_{\alpha/2}=-z_{1-\alpha/2}$،
\begin{align*}
1-\alpha
&=\Prob\left(-z_{1-\alpha/2}\le\frac{\bar X-\mu}{\sigma/\sqrt n}\le z_{1-\alpha/2}\right)\\
&=\Prob\left(\bar X-z_{1-\alpha/2}\frac\sigma{\sqrt n}
\le\mu\le\bar X+z_{1-\alpha/2}\frac\sigma{\sqrt n}\right).
\end{align*}
\begin{result}{فاصلهٔ اطمینان دقیق برای میانگین؛ $\sigma$ معلوم}
\[
\boxed{\bar x\ \pm\ z_{1-\alpha/2}\frac\sigma{\sqrt n}.}
\]
در این نمادگذاری، $c\pm d$ به معنی بازهٔ $[c-d,c+d]$ است. طول بازه برابر $2z_{1-\alpha/2}\sigma/\sqrt n$ خواهد بود.
\end{result}
\begin{example}{نمونهٔ نرمال با واریانس ۲۵}
برای $n=16$، $\bar x=35$ و $\sigma^2=25$، فاصلهٔ اطمینان ۹۵ درصد برای $\mu$ برابر است با
\[
35\pm1.95996\frac5{\sqrt{16}}
=35\pm2.44995
=\boxed{[32.550,\ 37.450]}.
\]
\end{example}
\subsection{چندک‌های رایج و طول فاصله}
\begin{center}
\begin{tabular}{c|ccc}
\toprule
سطح اطمینان&$90\%$&$95\%$&$99\%$\\
\midrule
$z_{1-\alpha/2}$&$1.64485$&$1.95996$&$2.57583$\\
\bottomrule
\end{tabular}
\end{center}
با ثابت بودن $n$ و $\sigma$، افزایش سطح اطمینان بازه را بلندتر می‌کند؛ با ثابت بودن سطح اطمینان، افزایش حجم نمونه بازه را کوتاه‌تر می‌کند.

می‌توان به‌جای دو دُم برابر، $a<b$ را چنان انتخاب کرد که $\Prob(a<Z<b)=1-\alpha$. در این صورت بازهٔ میانگین برابر است با
\[\left[\bar x-b\frac\sigma{\sqrt n},\ \bar x-a\frac\sigma{\sqrt n}\right].\]
برای چگالی نرمال متقارن و تک‌قله‌ای، بازهٔ مرکزی با دو دُم برابر کوتاه‌ترین بازه را در میان این بازه‌های با ضرایب ثابت و پوشش یکسان دارد.
\section{توزیع‌های نمونه‌ای کای‌دو و استیودنت}
\source{۱۲ و ۱۳}
اگر $Z_1,\ldots,Z_\nu$ \textbf{مستقل} و هر یک $\N(0,1)$ باشند،
\[Q=\sum_{i=1}^{\nu}Z_i^2\sim\chi^2_\nu.\]
به‌ویژه $Z_i^2\sim\chi^2_1$. اگر $Z\sim\N(0,1)$ و $Q\sim\chi^2_\nu$ مستقل باشند،
\[T=\frac{Z}{\sqrt{Q/\nu}}\sim t_\nu.\]
توزیع $t$ نسبت به صفر متقارن و دارای دُم‌های سنگین‌تر از نرمال استاندارد است؛ با افزایش درجات آزادی به نرمال استاندارد نزدیک می‌شود.
\begin{center}
\begin{LTR}
\begin{tikzpicture}[x=1.35cm,y=5cm]
\draw[muted,-{Latex}] (-3.9,0)--(4,0) node[right] {$x$};
\draw[navy,thick] plot[domain=-3.8:3.8,samples=120] (\x,{exp(-\x*\x/2)/sqrt(2*pi)});
\draw[teal,thick,dashed] plot[domain=-3.8:3.8,samples=120] (\x,{2/(sqrt(3)*pi)*(1+\x*\x/3)^(-2)});
\draw[navy,thick] (1.8,.39)--(2.35,.39) node[right] {$\N(0,1)$};
\draw[teal,thick,dashed] (1.8,.32)--(2.35,.32) node[right] {$t_3$};
\end{tikzpicture}
\end{LTR}
\figcap{مقایسهٔ نرمال استاندارد و استیودنت با سه درجهٔ آزادی؛ سنگینی دُم‌ها در نمونهٔ کوچک مهم است.}
\end{center}
\subsection{نتیجهٔ بنیادی برای نمونهٔ نرمال}
برای نمونهٔ مستقل و هم‌توزیع نرمال،
\[
\frac{(n-1)S^2}{\sigma^2}\sim\chi^2_{n-1},
\qquad \bar X\ \text{و}\ S^2\ \text{مستقل‌اند}.
\]
ازاین‌رو
\[
\frac{\bar X-\mu}{S/\sqrt n}
=\frac{(\bar X-\mu)/(\sigma/\sqrt n)}
{\sqrt{[(n-1)S^2/\sigma^2]/(n-1)}}
\sim t_{n-1}.
\]
\textbf{چرایی درجات آزادی.} همانی مجموع مربعات می‌دهد
\[
\frac{\sum_i(X_i-\mu)^2}{\sigma^2}
=\frac{\sum_i(X_i-\bar X)^2}{\sigma^2}
+\frac{n(\bar X-\mu)^2}{\sigma^2}.
\]
در مدل نرمال، مؤلفهٔ میانگین یک بُعد و باقیمانده‌ها $n-1$ بُعد متعامد دارند. یک تبدیل متعامد از بردار نرمال استاندارد، مؤلفه‌های مستقل نرمال استاندارد تولید می‌کند؛ پس مجموع مربعات باقیمانده‌ها کای‌دو با $n-1$ درجهٔ آزادی است و از میانگین استقلال دارد.
\edit{صرفِ کم کردن یک متغیر کای‌دو از متغیر کای‌دوی دیگر، توزیع تفاضل را تعیین نمی‌کند. استقلال و تجزیهٔ متعامد نرمال لازم است. همچنین درجهٔ آزادی آمارهٔ $t$ در اینجا $n-1$ است، نه $n$.}
\section{میانگین یک جامعهٔ نرمال؛ واریانس نامعلوم}
\begin{result}{فاصلهٔ اطمینان دقیق استیودنت}
برای $X_i\iid\N(\mu,\sigma^2)$، $n\ge2$ و $\sigma$ نامعلوم،
\[
\boxed{\bar x\ \pm\ t_{1-\alpha/2,n-1}\frac{s}{\sqrt n}.}
\]
این فرمول از حل نامساوی $-t_{1-\alpha/2,n-1}\le(\bar X-\mu)/(S/\sqrt n)\le t_{1-\alpha/2,n-1}$ به دست می‌آید.
\end{result}
\begin{example}{سرعت پنج خودرو}
سرعت پنج خودرو از جامعه‌ای نرمال اندازه‌گیری شده است؛ میانگین نمونه $80$ و انحراف معیار نمونه $6$ کیلومتر بر ساعت است. چون $n=5$ و $t_{0.975,4}=2.77645$، فاصلهٔ اطمینان ۹۵ درصد میانگین سرعت برابر است با
\[
80\pm2.77645\frac6{\sqrt5}
=\boxed{[72.550,\ 87.450]}.
\]
\end{example}
\section{حد خطا و تعیین حجم نمونه}
\source{۱۴}
نیم‌طول بازه را حد خطا یا حاشیهٔ خطا می‌نامیم. در حالت واریانس معلوم،
\[E=z_{1-\alpha/2}\frac\sigma{\sqrt n}.\]
رویداد $|\bar X-\mu|\le E$ احتمال $1-\alpha$ دارد؛ این کران، تضمین قطعی خطای مشاهده‌شده نیست.
\begin{result}{حجم نمونه برای حاشیهٔ خطای حداکثر $e$}
\[
E\le e\quad\Longleftrightarrow\quad
n\ge\left(\frac{z_{1-\alpha/2}\sigma}{e}\right)^2.
\]
حداقل حجم نمونه، کوچک‌ترین عدد صحیحِ بزرگ‌تر یا مساوی سمت راست است.
\end{result}
اگر $\sigma$ نامعلوم باشد، در مرحلهٔ برنامه‌ریزی از یک برآورد مقدماتی یا مطالعهٔ پیشین استفاده می‌کنیم. رابطهٔ $t_{1-\alpha/2,n-1}s/\sqrt n\le e$ به $n$ وابستگی ضمنی دارد و معمولاً به‌صورت تکراری حل می‌شود؛ نمی‌توان $s$ آینده را پیش از نمونه‌گیری دانست.
\begin{example}{حجم نمونه با انحراف معیار ۴۰}
اگر $\sigma=40$، سطح اطمینان ۹۵ درصد و حد خطای مطلوب $10$ باشد،
\[
n\ge\left(\frac{1.95996\times40}{10}\right)^2=61.4654\ldots
\quad\Longrightarrow\quad\boxed{n_{\min}=62}.
\]
\edit{در متن مبنا محاسبهٔ عددی ناتمام و مقدار میانی کم‌دقت بود؛ حجم نمونه باید رو به بالا گرد شود.}
\end{example}
\section{قضیهٔ حد مرکزی و فاصله‌های تقریبی}
\source{۱۵}
\begin{result}{قضیهٔ حد مرکزی برای نمونهٔ مستقل و هم‌توزیع}
اگر $\E(X_i)=\mu$ و $0<\Var(X_i)=\sigma^2<\infty$،
\[
\frac{\sqrt n(\bar X_n-\mu)}\sigma\xrightarrow{d}\N(0,1).
\]
در نتیجه برای هر $a\in\R$،
\[
\Prob\left(\frac{\bar X_n-\mu}{\sigma/\sqrt n}\le a\right)\to\Phi(a).
\]
\end{result}
برای نمونهٔ کافی بزرگ، $\bar X_n\approx\N(\mu,\sigma^2/n)$. اگر $S^2$ برآوردگر سازگار واریانس باشد، جایگزینی $\sigma$ با $S$ نیز به تقریب نرمال می‌انجامد. سرعت مناسب شدن تقریب به شکل توزیع، چولگی و دُم‌ها بستگی دارد؛ یک آستانهٔ ثابت برای همهٔ جوامع وجود ندارد.
\section{اختلاف میانگین دو جامعهٔ مستقل}
\source{۱۵ تا ۱۸}
دو نمونهٔ مستقل از هم در نظر بگیرید:
\[
X_{1i}\iid\N(\mu_1,\sigma_1^2),\quad i=1,\ldots,n_1,
\qquad X_{2j}\iid\N(\mu_2,\sigma_2^2),\quad j=1,\ldots,n_2.
\]
هدف، برآورد $\Delta=\mu_1-\mu_2$ است و
\[
\bar X_1-\bar X_2\sim\N\left(\Delta,\frac{\sigma_1^2}{n_1}+\frac{\sigma_2^2}{n_2}\right).
\]
\subsection{واریانس‌های معلوم}
\begin{result}{فاصلهٔ اطمینان اختلاف میانگین؛ حالت معلوم}
\[
\boxed{(\bar x_1-\bar x_2)\ \pm\ z_{1-\alpha/2}
\sqrt{\frac{\sigma_1^2}{n_1}+\frac{\sigma_2^2}{n_2}}.}
\]
کمیت محوری برابر است با تفاضلِ برآورد و پارامتر، تقسیم بر خطای استاندارد بالا؛ توزیع آن نرمال استاندارد است.
\end{result}
\begin{example}{مقایسهٔ میانگین نمرهٔ دو گروه}
دو گروه مستقل با $n_1=10,n_2=12$ و میانگین‌های $\bar x_1=14,\bar x_2=12$ داریم. با فرض واریانس‌های معلوم $\sigma_1^2=40,\sigma_2^2=60$، فاصلهٔ اطمینان ۹۰ درصد برای $\mu_1-\mu_2$ برابر است با
\[
2\pm1.64485\sqrt{40/10+60/12}
=2\pm4.93456=\boxed{[-2.935,\ 6.935]}.
\]
\edit{رقم واریانس دوم در صفحهٔ ۱۶ کم‌خواناست؛ در حل بالا «۶۰» در نظر گرفته شده است. اگر خوانش درست «۹۰» باشد، همان فرمول بازهٔ $[-3.578,\ 7.578]$ را می‌دهد.}
\end{example}
\subsection{واریانس‌های نامعلوم ولی برابر}
فرض کنید $\sigma_1^2=\sigma_2^2=\sigma^2$. برآوردگر تجمیعی واریانس برابر است با
\[
S_p^2=\frac{(n_1-1)S_1^2+(n_2-1)S_2^2}{n_1+n_2-2}.
\]
زیرا امید صورت برابر $(n_1+n_2-2)\sigma^2$ است، $\E(S_p^2)=\sigma^2$. همچنین از استقلال دو نمونه و جمع‌پذیری کای‌دوهای مستقل،
\[
\frac{(n_1+n_2-2)S_p^2}{\sigma^2}\sim\chi^2_{n_1+n_2-2}.
\]
این کمیت از اختلاف میانگین‌ها مستقل است. بنابراین
\[
T=\frac{(\bar X_1-\bar X_2)-\Delta}
{S_p\sqrt{1/n_1+1/n_2}}\sim t_{n_1+n_2-2}.
\]
\begin{result}{فاصلهٔ اطمینان تجمیعی}
\[
\boxed{(\bar x_1-\bar x_2)\ \pm\ t_{1-\alpha/2,n_1+n_2-2}
s_p\sqrt{\frac1{n_1}+\frac1{n_2}}.}
\]
فرض برابری واریانس‌ها بخشی از مدل این روش است.
\end{result}
\subsection{واریانس‌های نامعلوم و نامساوی؛ روش ولچ}
وقتی برابری واریانس‌ها را فرض نمی‌کنیم،
\[
T_W=\frac{(\bar X_1-\bar X_2)-\Delta}
{\sqrt{S_1^2/n_1+S_2^2/n_2}}
\approx t_\nu,
\]
که درجهٔ آزادی تقریبی ساترویت بر اساس داده‌ها برابر است با
\[
\nu=\frac{(s_1^2/n_1+s_2^2/n_2)^2}
{\dfrac{(s_1^2/n_1)^2}{n_1-1}+\dfrac{(s_2^2/n_2)^2}{n_2-1}}.
\]
\begin{result}{فاصلهٔ اطمینان تقریبی ولچ}
\[
\boxed{(\bar x_1-\bar x_2)\ \pm\ t_{1-\alpha/2,\nu}
\sqrt{\frac{s_1^2}{n_1}+\frac{s_2^2}{n_2}}.}
\]
$\nu$ می‌تواند غیرصحیح باشد. اگر جدول فقط درجات آزادی صحیح داشته باشد، گرد کردن رو به پایین معمولاً انتخاب محافظه‌کارانه‌تری است.
\end{result}
\edit{در دست‌نویس، توزیع ولچ گاه با علامت برابریِ توزیعی نوشته شده است؛ در این روش رابطه، تقریبی است.}
\section{فاصلهٔ اطمینان واریانس جامعهٔ نرمال}
\source{۱۹ و ۲۰}
\subsection{میانگین معلوم}
اگر $\mu$ معلوم باشد،
\[
Q=\frac{\sum_{i=1}^n(X_i-\mu)^2}{\sigma^2}\sim\chi^2_n.
\]
قرار دهید $A=\sum_i(x_i-\mu)^2$. از
\[\chi^2_{\alpha/2,n}\le A/\sigma^2\le\chi^2_{1-\alpha/2,n}\]
با معکوس‌کردن و حل نسبت به $\sigma^2$ به دست می‌آید
\[
\boxed{\left[\frac{A}{\chi^2_{1-\alpha/2,n}},\
\frac{A}{\chi^2_{\alpha/2,n}}\right].}
\]
\subsection{میانگین نامعلوم}
در این حالت، $A=(n-1)s^2$ و درجهٔ آزادی $n-1$ است:
\begin{result}{فاصلهٔ اطمینان دقیق واریانس؛ میانگین نامعلوم}
\[
\boxed{\left[\frac{(n-1)s^2}{\chi^2_{1-\alpha/2,n-1}},\
\frac{(n-1)s^2}{\chi^2_{\alpha/2,n-1}}\right].}
\]
برای فاصلهٔ اطمینان $\sigma$، از دو کران مثبت ریشهٔ دوم می‌گیریم.
\end{result}
\begin{center}
\begin{LTR}
\begin{tikzpicture}[x=.57cm,y=12cm]
\fill[teal!17] plot[domain=1.635:12.592,samples=100] (\x,{\x*\x*exp(-\x/2)/16})--(12.592,0)--(1.635,0)--cycle;
\draw[navy,thick] plot[domain=.01:19,samples=140] (\x,{\x*\x*exp(-\x/2)/16});
\draw[muted,-{Latex}] (0,0)--(19.5,0) node[right] {$q$};
\draw[gold,dashed] (1.635,0)--(1.635,{1.635^2*exp(-1.635/2)/16});
\draw[gold,dashed] (12.592,0)--(12.592,{12.592^2*exp(-12.592/2)/16});
\node[below] at (1.635,0) {$\chi^2_{\alpha/2,\nu}$};
\node[below] at (12.592,0) {$\chi^2_{1-\alpha/2,\nu}$};
\node[teal] at (6.4,.055) {$1-\alpha$};
\end{tikzpicture}
\end{LTR}
\figcap{بازهٔ مرکزی زیر چگالی کای‌دو؛ شکل برای $\nu=6$ و $\alpha=0.10$ رسم شده است. فاصلهٔ واریانس معمولاً متقارن نیست.}
\end{center}
\begin{example}{نمونهٔ ده‌تایی با واریانس ۵}
از جامعهٔ نرمال، نمونه‌ای با $n=10$ و $s^2=5$ گرفته‌ایم. برای اطمینان ۹۰ درصد،
\[
\chi^2_{0.95,9}\approx16.91898,\qquad
\chi^2_{0.05,9}\approx3.32511.
\]
بنابراین
\[
\left[\frac{45}{16.91898},\frac{45}{3.32511}\right]
=\boxed{[2.660,\ 13.533]}.
\]
\end{example}
\section{نسبت واریانس دو جامعهٔ نرمال}
\source{۲۰ تا ۲۲}
\begin{definition}{توزیع $F$}
اگر $U\sim\chi^2_{\nu_1}$ و $V\sim\chi^2_{\nu_2}$ مستقل باشند،
\[
F=\frac{U/\nu_1}{V/\nu_2}\sim F_{\nu_1,\nu_2}.
\]
ترتیب دو درجهٔ آزادی مهم است: اولی به صورت و دومی به مخرج مربوط می‌شود.
\end{definition}
اگر $F\sim F_{\nu_1,\nu_2}$، آنگاه $1/F\sim F_{\nu_2,\nu_1}$ و
\[
\boxed{\frac1{F_{p;\nu_1,\nu_2}}=F_{1-p;\nu_2,\nu_1}.}
\]
\textbf{اثبات چندک معکوس.} چون تبدیل معکوس نزولی و توزیع پیوسته است،
\[
\Prob\left(\frac1F\le\frac1{F_{p;\nu_1,\nu_2}}\right)
=\Prob(F\ge F_{p;\nu_1,\nu_2})=1-p.
\]
\subsection{یک فرمول مشترک برای هر دو حالت میانگین}
برای دو نمونهٔ مستقل نرمال تعریف کنید
\[
R_j=\begin{cases}
\displaystyle\frac1{n_j}\sum_i(X_{ji}-\mu_j)^2,&\mu_j\text{ معلوم،}\quad\nu_j=n_j,\\[4pt]
S_j^2,&\mu_j\text{ نامعلوم،}\quad\nu_j=n_j-1.
\end{cases}
\]
در هر دو حالت، برای $\rho=\sigma_1^2/\sigma_2^2$،
\[Q=\frac{R_1/R_2}{\rho}\sim F_{\nu_1,\nu_2}.\]
\begin{result}{فاصلهٔ اطمینان نسبت واریانس‌ها}
با $r_j$ به‌عنوان مقدار مشاهده‌شدهٔ $R_j$،
\[
\boxed{\left[\frac{r_1/r_2}{F_{1-\alpha/2;\nu_1,\nu_2}},\
\frac{r_1/r_2}{F_{\alpha/2;\nu_1,\nu_2}}\right].}
\]
فرم معادل با چندک‌های درجات آزادی معکوس عبارت است از
\[
\left[\frac{r_1}{r_2}F_{\alpha/2;\nu_2,\nu_1},\
\frac{r_1}{r_2}F_{1-\alpha/2;\nu_2,\nu_1}\right].
\]
\end{result}
اگر فاصله شامل $1$ باشد، آزمون دوطرفهٔ متناظر برابری واریانس‌ها را رد نمی‌کند؛ این نتیجه اثبات برابری نیست.
\begin{example}{دو نمونه با واریانس‌های ۸ و ۱۶}
در صفحهٔ ۲۲، نمونهٔ اول $n_1=10,\bar x_1=10,s_1^2=8$ و نمونهٔ دوم $n_2=6,\bar x_2=8,s_2^2=16$ دارد. نسبت نمونه‌ای برابر $1/2$ است. فاصلهٔ اطمینان ۹۵ درصد برای $\rho$،
\[
\left[\frac{0.5}{F_{0.975;9,5}},\frac{0.5}{F_{0.025;9,5}}\right]
=\boxed{[0.07484,\ 2.24221]}.
\]
در سطح ۵ درصد، شواهد کافی برای رد برابری واریانس‌ها نداریم.

\textbf{مقایسهٔ میانگین‌ها در همان مثال.} متن مبنا دربارهٔ برابری میانگین‌ها نیز پرسیده است؛ این پرسش با فاصلهٔ نسبت واریانس پاسخ داده نمی‌شود. با روش ولچ و $\nu\approx8.048$، فاصلهٔ ۹۵ درصد اختلاف میانگین‌ها تقریباً $[-2.289,6.289]$ است و صفر را دربرمی‌گیرد.
\edit{میانگین و واریانس دو پارامتر متفاوت‌اند. برای بررسی برابری میانگین‌ها باید فاصله یا آزمون اختلاف میانگین را به کار برد.}
\end{example}
\section{داده‌های زوجی و اختلاف میانگین‌ها}
\source{۲۳ و ۲۴}
در داده‌های زوجی، $X_i$ و $Y_i$ به یک واحد یا یک زوج مرتبط مربوط‌اند. وابستگی داخل هر زوج مجاز است؛ فرض می‌کنیم زوج‌های مختلف مستقل و هم‌توزیع‌اند. اختلاف‌ها را تعریف کنید:
\[
D_i=X_i-Y_i,\qquad
\mu_D=\mu_X-\mu_Y,
\qquad\sigma_D^2=\sigma_X^2+\sigma_Y^2-2\Cov(X_i,Y_i).
\]
اگر اختلاف‌ها نرمال باشند،
\[
\bar D=\frac1n\sum_iD_i=\bar X-\bar Y,\qquad
S_D^2=\frac1{n-1}\sum_i(D_i-\bar D)^2,
\qquad
\frac{\bar D-\mu_D}{S_D/\sqrt n}\sim t_{n-1}.
\]
\begin{result}{فاصلهٔ اطمینان زوجی}
\[\boxed{\bar d\ \pm\ t_{1-\alpha/2,n-1}\frac{s_D}{\sqrt n}.}\]
داده‌های زوجی را با فرمول دو نمونهٔ مستقل تحلیل نکنید؛ کوواریانس داخل زوج بخشی از اطلاعات مسئله است.
\end{result}
\edit{نرمال بودن جداگانهٔ $X_i$ و $Y_i$ برای نرمال بودن اختلاف کافی نیست. فرض مستقیمِ نرمال بودن اختلاف‌ها یا نرمال مشترک بودن هر زوج، شرط درست است.}
\begin{example}{نمرات روش‌های آماری و جبر خطی}
نمرات پنج دانشجو در دو درس به شکل زیر است؛ جهت اختلاف، «روش‌های آماری منهای جبر خطی» انتخاب شده است.
\begin{center}
\begin{tabular}{r|ccccc}
\toprule
دانشجو&$1$&$2$&$3$&$4$&$5$\\
\midrule
روش‌های آماری&$17$&$15$&$12$&$18$&$14$\\
جبر خطی&$19$&$14$&$10$&$13$&$16$\\
$D_i$&$-2$&$1$&$2$&$5$&$-2$\\
\bottomrule
\end{tabular}
\end{center}
\[
\bar d=\frac45=0.8,\qquad \sum_i d_i^2=38,
\qquad s_D^2=\frac{38-5(0.8)^2}{4}=\boxed{8.7}.
\]
با فرض نرمال بودن اختلاف‌ها، فاصلهٔ ۹۵ درصد برابر است با
\[
0.8\pm2.77645\sqrt{\frac{8.7}{5}}
=\boxed{[-2.862,\ 4.462]}.
\]
وجود صفر در فاصله به معنی نبود شواهد کافی برای اختلاف میانگین دو درس در این سطح است، نه اثبات برابر بودن آن‌ها.
\edit{در دست‌نویس، پس از نوشتن فرمول صحیح واریانس، جملهٔ $n\bar d^2$ در محاسبه حذف شده و عدد $9.5$ به دست آمده است. مقدار درست $8.7$ و بازهٔ درست همان بازهٔ بالا است.}
\end{example}
\section{فاصلهٔ اطمینان نسبت یک جامعه}
\source{۲۴ و ۲۵}
اگر $X_i\iid\operatorname{Bern}(p)$، برآوردگر نسبت جامعه برابر است با
\[
\widehat p=\frac1n\sum_iX_i,\qquad
\E(\widehat p)=p,\qquad \Var(\widehat p)=\frac{p(1-p)}n.
\]
برای $0<p<1$ و نمونهٔ بزرگ،
\[
\frac{\widehat p-p}{\sqrt{p(1-p)/n}}\xrightarrow{d}\N(0,1).
\]
با سازگاری $\widehat p$ و جایگزینی پارامتر در خطای استاندارد،
\[
\frac{\widehat p-p}{\sqrt{\widehat p(1-\widehat p)/n}}
\xrightarrow{d}\N(0,1).
\]
\begin{result}{فاصلهٔ تقریبی والد برای نسبت}
\[
\boxed{\widehat p\ \pm\ z_{1-\alpha/2}
\sqrt{\frac{\widehat p(1-\widehat p)}n}.}
\]
این فاصله دقیق نیست. وقتی تعداد موفقیت‌ها یا شکست‌ها کم باشد، تقریب ممکن است پوشش نامناسبی داشته باشد؛ در $\widehat p=0$ یا $1$ نیز به بازهٔ تک‌نقطه‌ای ناموجه تبدیل می‌شود.
\end{result}
\begin{example}{پنج قطعهٔ معیوب از صد قطعه}
اگر در نمونه‌ای با حجم $100$، تعداد $5$ قطعه معیوب باشد، $\widehat p=0.05$. بازهٔ والد با اطمینان اسمی ۹۵ درصد برابر است با
\[
0.05\pm1.95996\sqrt{\frac{0.05(0.95)}{100}}
=\boxed{[0.00728,\ 0.09272]}.
\]
\textbf{نکتهٔ کاربردی.} تعداد عیب‌ها فقط پنج است؛ این مثال در مرز یا خارج از شرایط مناسب تقریب نرمال قرار می‌گیرد. بازهٔ بالا پاسخ روش معرفی‌شده در جزوه است، ولی برای گزارش عملی بهتر است روش‌هایی مانند ویلسون یا روش دقیق دوجمله‌ای بررسی شوند.
\end{example}
\section{فاصلهٔ اطمینان اختلاف دو نسبت}
\source{۲۶}
برای دو نمونهٔ برنولی مستقل با اندازه‌های $n_1,n_2$ و نسبت‌های $p_1,p_2$،
\[
\E(\widehat p_1-\widehat p_2)=p_1-p_2,
\qquad
\Var(\widehat p_1-\widehat p_2)=\frac{p_1(1-p_1)}{n_1}+\frac{p_2(1-p_2)}{n_2}.
\]
با بزرگ بودن نمونه‌ها و مناسب بودن تقریب،
\begin{result}{فاصلهٔ تقریبی اختلاف دو نسبت مستقل}
\[
\boxed{(\widehat p_1-\widehat p_2)\ \pm\ z_{1-\alpha/2}
\sqrt{\frac{\widehat p_1(1-\widehat p_1)}{n_1}
+\frac{\widehat p_2(1-\widehat p_2)}{n_2}}.}
\]
برای ساخت این فاصله، نسبت‌ها را جداگانه در خطای استاندارد جایگزین می‌کنیم؛ نسبت تجمیعیِ آزمون برابری نسبت‌ها در این فرمول به کار نمی‌رود.
\end{result}
\begin{example}{نظر دانشجویان دختر و پسر دربارهٔ یک برنامه}
از $200$ دانشجوی دختر و $200$ دانشجوی پسر، به‌ترتیب $128$ و $72$ نفر موافق برنامه‌اند. با انتخاب جهت «دختر منهای پسر»،
\[
\widehat p_1=0.64,\qquad\widehat p_2=0.36,\qquad
\widehat p_1-\widehat p_2=0.28.
\]
فاصلهٔ اطمینان ۹۵ درصد برابر است با
\[
0.28\pm1.95996\sqrt{\frac{0.64(0.36)}{200}+\frac{0.36(0.64)}{200}}
=\boxed{[0.18592,\ 0.37408]}.
\]
با این جهت‌گذاری، نسبت موافقت دختران حدود ۱۸٫۶ تا ۳۷٫۴ واحد درصد بیشتر برآورد می‌شود. برای «پسر منهای دختر»، علامت هر دو کران و ترتیب آن‌ها معکوس می‌شود.
\end{example}

\chapter{آزمون فرض آماری}
\source{۲۷ تا ۴۵}
\section{فرض آماری، آماره و ناحیهٔ رد}
\begin{definition}{اجزای یک مسئلهٔ آزمون فرض}
\textbf{فرض آماری} گزاره‌ای دربارهٔ توزیع جامعه یا پارامترهای آن است. فرض مبنا را با $H_0$ و فرض مقابل را با $H_1$ نشان می‌دهیم.

\textbf{فرض ساده} یک توزیع را به‌طور کامل مشخص می‌کند. فرضی که بیش از یک توزیع را مجاز بداند، مرکب است؛ حتی تثبیت یک پارامتر در حضور پارامتر مزاحم نامعلوم می‌تواند فرضی مرکب بسازد.

\textbf{آمارهٔ آزمون} تابعی از نمونه است که بر اساس مقدار آن دربارهٔ رد یا رد نکردن $H_0$ تصمیم می‌گیریم.

\textbf{ناحیهٔ رد} یا بحرانی، مجموعهٔ مقادیر آماره یا نمونه است که به رد $H_0$ منجر می‌شود؛ آن را با $C$ نشان می‌دهیم.
\end{definition}
برای مثال، در آزمون میانگین، سه فرض مقابل رایج عبارت‌اند از
\[
H_1:\mu<\mu_0,\qquad H_1:\mu>\mu_0,\qquad H_1:\mu\ne\mu_0.
\]
این‌ها به‌ترتیب آزمون چپ‌طرفه، راست‌طرفه و دوطرفه را تعیین می‌کنند. در هر سه حالت، مرز فرض صفر $\mu=\mu_0$ است. جهت فرض مقابل باید بر اساس پرسش علمی و پیش از دیدن داده‌ها تعیین شود.
\begin{important}{رد نکردن، اثبات نیست}
تصمیم معمول آزمون «رد $H_0$» یا «رد نکردن $H_0$» است. رد نکردن ممکن است از کم بودن شواهد یا توان پایین ناشی شود و به معنای اثبات $H_0$ نیست. در این نسخه، عبارت «قبول فرض صفر» در متن مبنا با همین تعبیر دقیق‌تر جایگزین شده است.
\end{important}
\section{خطاهای نوع اول و دوم؛ توان آزمون}
\source{۲۷ و ۲۸}
\begin{center}
\begin{tabular}{r|cc}
\toprule
تصمیم & $H_0$ درست است & $H_0$ نادرست است\\
\midrule
رد $H_0$ & خطای نوع اول & تصمیم درست\\
رد نکردن $H_0$ & تصمیم درست & خطای نوع دوم\\
\bottomrule
\end{tabular}
\end{center}
\begin{result}{سطح و توان}
برای یک فرض صفر ساده،
\[
\alpha=\Prob_{H_0}(\text{رد }H_0).
\]
برای یک مقدار مشخص $\theta\in\Theta_1$ از فرض مقابل،
\[
\beta(\theta)=\Prob_\theta(\text{رد نکردن }H_0),\qquad
\pi(\theta)=1-\beta(\theta)=\Prob_\theta(\text{رد }H_0).
\]
تابع $\pi(\theta)$ تابع توان است. برای فرض صفر مرکب، سطح آزمون با بیشترین احتمال خطای نوع اول روی مقادیر مجاز فرض صفر کنترل می‌شود.
\end{result}
معمولاً $\alpha$ را پیشاپیش در سطحی مانند $0.05$ تثبیت می‌کنیم و میان آزمون‌های دارای این کنترل، توان بیشتر را مطلوب می‌دانیم. با حجم نمونهٔ ثابت، کاهش $\alpha$ اغلب با افزایش $\beta$ همراه است؛ بزرگ‌تر کردن نمونه می‌تواند هر دو خطا را بهبود دهد.
\begin{center}
\begin{LTR}
\begin{tikzpicture}[x=1.2cm,y=4.1cm]
\fill[alert!20] plot[domain=1.645:4,samples=70] (\x,{exp(-\x*\x/2)/sqrt(2*pi)})--(4,0)--(1.645,0)--cycle;
\fill[teal!17] plot[domain=-1:1.645,samples=70] (\x,{exp(-(\x-2.5)^2/2)/sqrt(2*pi)})--(1.645,0)--(-1,0)--cycle;
\draw[navy,thick] plot[domain=-3:5.5,samples=120] (\x,{exp(-\x*\x/2)/sqrt(2*pi)});
\draw[teal,thick] plot[domain=-3:5.5,samples=120] (\x,{exp(-(\x-2.5)^2/2)/sqrt(2*pi)});
\draw[muted,-{Latex}] (-3.1,0)--(5.8,0);
\draw[gold,dashed] (1.645,0)--(1.645,.44);
\node[navy,above] at (0,.4) {$H_0$};\node[teal,above] at (2.5,.4) {$H_1$};
\node[below] at (1.645,0) {$c$};
\node[alert] at (2.15,.035) {$\alpha$};\node[teal] at (.95,.06) {$\beta$};
\end{tikzpicture}
\end{LTR}
\figcap{نمایش خطای نوع اول و دوم برای یک آزمون راست‌طرفه در برابر یک مقدار مشخص از فرض مقابل. ناحیهٔ رد، سمت راست $c$ است.}
\end{center}
\begin{example}{ده پرتاب سکه}
فرض کنید $X_i$ شاخص شیر آمدن در پرتاب $i$ باشد و
\[
X_i\iid\operatorname{Bern}(p),\quad n=10,\qquad
H_0:p=\frac12,\quad H_1:p=\frac34.
\]
ناحیهٔ رد را $C=\{\sum_iX_i\ge7\}$ انتخاب می‌کنیم. اگر $Y=\sum_iX_i$، آنگاه
\begin{align*}
\alpha&=\Prob_{p=1/2}(Y\ge7)
=\sum_{y=7}^{10}\binom{10}{y}2^{-10}
=\boxed{0.171875},\\
\beta&=\Prob_{p=3/4}(Y\le6)
=\sum_{y=0}^{6}\binom{10}{y}(3/4)^y(1/4)^{10-y}
\approx\boxed{0.224125}.
\end{align*}
توان در $p=3/4$ برابر $1-\beta\approx0.775875$ است. این آزمونِ ازپیش‌تعریف‌شده، سطح ۵ درصد ندارد؛ سطح واقعی آن همان $0.171875$ است.
\end{example}
\section{آزمون میانگین یک جامعهٔ نرمال؛ واریانس معلوم}
\source{۲۹ تا ۳۱}
فرض کنید $X_i\iid\N(\mu,\sigma^2)$ و $\sigma$ معلوم باشد. برای $H_0:\mu=\mu_0$، آمارهٔ
\[Z=\frac{\bar X-\mu_0}{\sigma/\sqrt n}\]
تحت فرض صفر، نرمال استاندارد است.
\subsection{استخراج ناحیهٔ رد چپ‌طرفه}
برای $H_1:\mu<\mu_0$، میانگین‌های کوچک به سود فرض مقابل‌اند. ناحیهٔ رد را $\bar X-\mu_0<c$ می‌گیریم. از
\[
\alpha=\Prob_{H_0}(\bar X-\mu_0<c)
=\Prob\left(Z<\frac{c}{\sigma/\sqrt n}\right)
\]
نتیجه می‌شود $c=z_\alpha\sigma/\sqrt n=-z_{1-\alpha}\sigma/\sqrt n$. پس ناحیهٔ رد $Z<-z_{1-\alpha}$ است.
\begin{result}{سه آزمون میانگین با واریانس معلوم؛ سطح $\alpha$}
با مقدار مشاهده‌شدهٔ $z_{\rm obs}$،
\[
\begin{array}{c|c}
\text{فرض مقابل}&\text{شرط رد فرض صفر}\\\hline
\mu<\mu_0&z_{\rm obs}<-z_{1-\alpha}\\[2pt]
\mu>\mu_0&z_{\rm obs}>z_{1-\alpha}\\[2pt]
\mu\ne\mu_0&|z_{\rm obs}|>z_{1-\alpha/2}
\end{array}
\]
\end{result}
برای آزمون راست‌طرفه، از $\Prob(Z<c/(\sigma/\sqrt n))=1-\alpha$ نتیجه می‌شود $c=z_{1-\alpha}\sigma/\sqrt n$. برای آزمون دوطرفه،
\[
\alpha=\Prob(|Z|>c/(\sigma/\sqrt n))
=2\left[1-\Phi\left(\frac c{\sigma/\sqrt n}\right)\right],
\]
بنابراین $c=z_{1-\alpha/2}\sigma/\sqrt n$.
\begin{example}{آزمون چپ‌طرفه}
داده‌ها عبارت‌اند از $n=16$، $\bar x=9$ و $\sigma^2=9$. فرض‌ها را
\[H_0:\mu=10,\qquad H_1:\mu<10\]
و سطح را $\alpha=0.05$ بگیرید. آنگاه
\[
z_{\rm obs}=\frac{9-10}{3/4}=-\frac43\approx-1.333.
\]
چون $-1.333>-1.64485$، فرض صفر رد نمی‌شود. مقدار $p$ حدود $0.09121$ است و همین تصمیم را تأیید می‌کند.
\end{example}
\begin{example}{آزمون دوطرفه با مقدار آمارهٔ ۲}
در مثال صفحهٔ ۳۱، $\bar x=3,\mu_0=1,\sigma^2=25$ و محاسبهٔ نوشته‌شده به $z_{\rm obs}=2$ می‌رسد. با خطای استاندارد $1$، حجم نمونهٔ متناظر $n=25$ است. برای
\[H_0:\mu=1,\qquad H_1:\mu\ne1,\qquad \alpha=0.05,\]
داریم $|2|>1.95996$؛ پس فرض صفر رد می‌شود. مقدار $p$ دوطرفه تقریباً $0.04550$ است.
\edit{زیرنویس حجم نمونه در صورت دست‌نویس کم‌خواناست؛ در این مثال، $n=25$ از خطای استاندارد و مقدار آمارهٔ همان حل بازسازی شده است.}
\end{example}
\section{آزمون میانگین یک جامعهٔ نرمال؛ واریانس نامعلوم}
\source{۳۲}
وقتی $\sigma$ نامعلوم باشد، آمارهٔ مناسب
\[T=\frac{\bar X-\mu_0}{S/\sqrt n}\sim t_{n-1}\quad\text{تحت }H_0\]
است. استخراج ناحیهٔ رد مانند حالت قبل انجام می‌شود و چندک‌های استیودنت جایگزین چندک‌های نرمال می‌شوند.
\begin{result}{سه آزمون استیودنت تک‌نمونه‌ای}
\[
\begin{array}{c|c}
H_1&\text{شرط رد }H_0\\\hline
\mu<\mu_0&t_{\rm obs}<-t_{1-\alpha,n-1}\\[2pt]
\mu>\mu_0&t_{\rm obs}>t_{1-\alpha,n-1}\\[2pt]
\mu\ne\mu_0&|t_{\rm obs}|>t_{1-\alpha/2,n-1}
\end{array}
\]
\end{result}
\begin{example}{آیا میانگین نمرات ریاضی برابر ۱۵ است؟}
در نمونه‌ای با حجم $9$، میانگین و انحراف معیار نمرات به‌ترتیب $13$ و $3$ است. با فرض نرمال بودن جامعه، برای
\[H_0:\mu=15,\qquad H_1:\mu\ne15\]
داریم
\[
t_{\rm obs}=\frac{13-15}{3/\sqrt9}=-2,\qquad
\nu=8,\qquad t_{0.975,8}=2.30600.
\]
چون $|-2|=2<2.30600$، در سطح ۵ درصد فرض صفر \textbf{رد نمی‌شود}. مقدار $p$ دوطرفه حدود $0.08052$ است.
\edit{در دست‌نویس مقدار بحرانی $1.86$، مربوط به $t_{0.95,8}$، برای آزمون دوطرفهٔ ۵ درصد استفاده و نتیجهٔ رد گرفته شده است. چندک درست $t_{0.975,8}$ و تصمیم درست، رد نکردن است.}
\end{example}
\section{آزمون اختلاف میانگین دو جامعهٔ مستقل}
\source{۳۳ و ۳۴}
در دست‌نویس، بعضی روابط با جهت $\mu_2-\mu_1$ نوشته شده‌اند. برای هماهنگی با فصل قبل، در این نسخه
\[\Delta=\mu_1-\mu_2,\qquad H_0:\Delta=\Delta_0\]
را به کار می‌بریم. تعویض جهت اختلاف، علامت آماره و جهت فرض‌های یک‌طرفه را هم عوض می‌کند.
\subsection{واریانس‌های معلوم}
آمارهٔ آزمون تحت فرض صفر برابر است با
\[
Z=\frac{(\bar X_1-\bar X_2)-\Delta_0}
{\sqrt{\sigma_1^2/n_1+\sigma_2^2/n_2}}\sim\N(0,1).
\]
برای $H_1:\Delta<\Delta_0$، $Z<-z_{1-\alpha}$؛ برای $H_1:\Delta>\Delta_0$، $Z>z_{1-\alpha}$؛ و برای $H_1:\Delta\ne\Delta_0$، $|Z|>z_{1-\alpha/2}$ شرط رد است.
\subsection{واریانس‌های نامعلوم ولی برابر}
از $S_p^2$ فصل قبل استفاده می‌کنیم:
\[
T_P=\frac{(\bar X_1-\bar X_2)-\Delta_0}
{S_p\sqrt{1/n_1+1/n_2}}\sim t_{n_1+n_2-2}.
\]
برای فرض مقابل کوچک‌تر، بزرگ‌تر و نابرابر، شرط‌های رد به‌ترتیب
\[
T_P<-t_{1-\alpha,\nu},\qquad
T_P>t_{1-\alpha,\nu},\qquad
|T_P|>t_{1-\alpha/2,\nu},\qquad \nu=n_1+n_2-2
\]
هستند.
\subsection{واریانس‌های نامعلوم و نامساوی}
\[
T_W=\frac{(\bar X_1-\bar X_2)-\Delta_0}
{\sqrt{S_1^2/n_1+S_2^2/n_2}}\approx t_\nu,
\]
که $\nu$ همان درجهٔ آزادی ساترویت در فصل قبل است. شرط‌های رد دقیقاً شکل سه شرط بالا را دارند، اما با $T_W$ و درجهٔ آزادی ولچ؛ کنترل سطح در این حالت تقریبی است.
\section{آزمون واریانس یک جامعهٔ نرمال}
\source{۳۵ و ۳۶}
فرض صفر $H_0:\sigma^2=\sigma_0^2$ است. بسته به معلوم بودن میانگین، دو آماره داریم:
\[
Q=\begin{cases}
\displaystyle\frac{\sum_i(X_i-\mu)^2}{\sigma_0^2},&\mu\text{ معلوم،}\quad\nu=n,\\[5pt]
\displaystyle\frac{(n-1)S^2}{\sigma_0^2},&\mu\text{ نامعلوم،}\quad\nu=n-1.
\end{cases}
\]
تحت فرض صفر، $Q\sim\chi^2_\nu$.

\textbf{استخراج برای فرض مقابل کوچک‌تر.} اگر $\mu$ معلوم باشد، قرار دهید $S_0^2=n^{-1}\sum_i(X_i-\mu)^2$. شرط رد را $S_0^2/\sigma_0^2<c$ بگیرید. از
\[
\alpha=\Prob_{H_0}\left(\frac{S_0^2}{\sigma_0^2}<c\right)
=\Prob(\chi^2_n<nc)
\]
نتیجه می‌شود $c=\chi^2_{\alpha,n}/n$. در حالت میانگین نامعلوم، جایگزینی $n$ با $n-1$ و $S_0^2$ با $S^2$ نتیجهٔ متناظر را می‌دهد.
\begin{result}{سه آزمون واریانس}
\[
\begin{array}{c|c}
H_1&\text{شرط رد }H_0\\\hline
\sigma^2<\sigma_0^2&Q<\chi^2_{\alpha,\nu}\\[2pt]
\sigma^2>\sigma_0^2&Q>\chi^2_{1-\alpha,\nu}\\[2pt]
\sigma^2\ne\sigma_0^2&Q<\chi^2_{\alpha/2,\nu}\ \text{یا}\ Q>\chi^2_{1-\alpha/2,\nu}
\end{array}
\]
در آزمون دوطرفهٔ این جزوه، احتمال هر دُم برابر $\alpha/2$ انتخاب شده است؛ چون کای‌دو نامتقارن است، شرط رد را با قدرمطلق یک حد متقارن نمی‌نویسیم.
\end{result}
\section{آزمون برابری واریانس دو جامعهٔ نرمال}
\source{۳۷ و ۳۸}
دو نمونهٔ نرمال مستقل در نظر بگیرید. از تعریف $R_j$ در فصل قبل استفاده می‌کنیم: اگر میانگین جامعه معلوم باشد، $R_j=n_j^{-1}\sum_i(X_{ji}-\mu_j)^2$ و $\nu_j=n_j$؛ در غیر این صورت $R_j=S_j^2$ و $\nu_j=n_j-1$.

برای فرض صفر $H_0:\sigma_1^2/\sigma_2^2=1$، آمارهٔ
\[F=\frac{R_1}{R_2}\sim F_{\nu_1,\nu_2}\]
تحت فرض صفر معتبر است. چون نسبت کوچک‌تر به سود واریانس اول کوچک‌تر است،
\begin{result}{سه آزمون نسبت واریانس‌ها}
\[
\begin{array}{c|c}
H_1&\text{شرط رد }H_0\\\hline
\sigma_1^2/\sigma_2^2<1&F<F_{\alpha;\nu_1,\nu_2}\\[2pt]
\sigma_1^2/\sigma_2^2>1&F>F_{1-\alpha;\nu_1,\nu_2}\\[2pt]
\sigma_1^2/\sigma_2^2\ne1&F<F_{\alpha/2;\nu_1,\nu_2}\ \text{یا}\ F>F_{1-\alpha/2;\nu_1,\nu_2}
\end{array}
\]
\end{result}
اگر فرض صفر $\sigma_1^2/\sigma_2^2=\rho_0$ باشد، آمارهٔ $F=(R_1/R_2)/\rho_0$ را با همین چندک‌ها مقایسه می‌کنیم.
\begin{important}{محدودیت مدل}
روابط دقیق کای‌دو و $F$ در این فصل به نرمال بودن جامعه‌ها متکی‌اند. صرف بزرگ بودن نمونه، فرض‌های دقیق این آزمون‌ها را جایگزین نمی‌کند. ترتیب صورت، مخرج و درجات آزادی را ثابت نگه دارید.
\end{important}
\section{آزمون میانگین در داده‌های زوجی}
\source{۳۹}
با تعریف $D_i=X_i-Y_i$، آزمون اختلاف میانگین‌ها به آزمون تک‌نمونه‌ای اختلاف‌ها تبدیل می‌شود. برای
\[H_0:\mu_D=\Delta_0\]
آمارهٔ
\[T_D=\frac{\bar D-\Delta_0}{S_D/\sqrt n}\sim t_{n-1}\]
را به کار می‌بریم؛ این رابطه در صورت نرمال بودن اختلاف‌ها دقیق است. در آزمون برابری میانگین‌ها $\Delta_0=0$ است. سه ناحیهٔ رد عبارت‌اند از
\[
\begin{array}{c|c}
H_1&\text{شرط رد }H_0\\\hline
\mu_D<\Delta_0&T_D<-t_{1-\alpha,n-1}\\[2pt]
\mu_D>\Delta_0&T_D>t_{1-\alpha,n-1}\\[2pt]
\mu_D\ne\Delta_0&|T_D|>t_{1-\alpha/2,n-1}
\end{array}
\]
استقلال لازم، استقلال \textbf{زوج‌ها از یکدیگر} است؛ استقلال دو عضو داخل هر زوج لازم نیست.
\section{آزمون نسبت یک جامعه}
\source{۴۰}
برای $X_i\iid\operatorname{Bern}(p)$ و $H_0:p=p_0$، با $0<p_0<1$، آمارهٔ آزمون بزرگ‌نمونه‌ای
\[Z=\frac{\widehat p-p_0}{\sqrt{p_0(1-p_0)/n}}\approx\N(0,1)\quad\text{تحت }H_0\]
است. برای فرض مقابل‌های $p<p_0$، $p>p_0$ و $p\ne p_0$ به‌ترتیب شرط‌های
\[
Z<-z_{1-\alpha},\qquad Z>z_{1-\alpha},\qquad |Z|>z_{1-\alpha/2}
\]
به کار می‌روند.
\begin{important}{در آزمون، واریانس را تحت فرض صفر بسازید}
برای این آزمون، خطای استاندارد شامل $p_0$ است؛ در فاصلهٔ والد از $\widehat p$ استفاده کردیم. شمار موفقیت و شکستِ مورد انتظار تحت فرض صفر باید برای تقریب نرمال کافی باشد. در نمونهٔ کوچک، آزمون دقیق دوجمله‌ای گزینهٔ مناسب‌تری است.
\end{important}
\section{آزمون برابری دو نسبت مستقل}
\source{۴۱ و ۴۲}
فرض کنید $X_{1i}\iid\operatorname{Bern}(p_1)$ و $X_{2j}\iid\operatorname{Bern}(p_2)$ و دو نمونه مستقل‌اند. هدف، آزمون
\[H_0:p_1=p_2\]
است. تحت فرض صفر، نسبت مشترک $p$ داریم و برآوردگر تجمیعی آن برابر است با
\[
\widehat p=\frac{\sum_iX_{1i}+\sum_jX_{2j}}{n_1+n_2}
=\frac{n_1\widehat p_1+n_2\widehat p_2}{n_1+n_2}.
\]
بنابراین آمارهٔ بزرگ‌نمونه‌ای آزمون برابر است با
\[
\boxed{Z=\frac{\widehat p_1-\widehat p_2}
{\sqrt{\widehat p(1-\widehat p)(1/n_1+1/n_2)}}.}
\]
برای فرض‌های مقابل $p_1<p_2$، $p_1>p_2$ و $p_1\ne p_2$، شرط‌های رد به‌ترتیب
\[
Z<-z_{1-\alpha},\qquad Z>z_{1-\alpha},\qquad |Z|>z_{1-\alpha/2}
\]
هستند. نسبت تجمیعی به این دلیل معتبر است که فرض صفر، برابری دو نسبت را تحمیل می‌کند. برای آزمون اختلاف غیرصفر $p_1-p_2=\Delta_0$، این تجمیع ساده عموماً معتبر نیست.
\section{مقدار احتمالی یا مقدار \lr{p}}
\source{۴۳ تا ۴۵}
\begin{definition}{مقدار احتمالی آزمون}
مقدار $p$ کوچک‌ترین سطح معنی‌داری در خانوادهٔ آزمون‌های مورد نظر است که در آن، دادهٔ مشاهده‌شده به رد $H_0$ منجر می‌شود. به‌طور معادل، تحت فرض صفر، احتمال مشاهدهٔ مقداری از آماره به اندازهٔ مقدار مشاهده‌شده یا افراطی‌تر در جهت فرض مقابل است.
\end{definition}
\begin{important}{مقدار $p$ چه چیزی را نمی‌گوید؟}
مقدار $p$ احتمال درست بودن $H_0$ نیست. این مقدار به مدل، فرض‌های آزمون و جهت فرض مقابل وابسته است؛ کوچک بودن آن نیز به‌تنهایی اندازه یا اهمیت عملی اثر را بیان نمی‌کند.
\end{important}
اگر آمارهٔ $T$ تحت فرض صفر تابع توزیع پیوستهٔ $G_0$ داشته باشد و مقدار مشاهده‌شده $t_{\rm obs}$ باشد،
\[
\begin{array}{c|c}
\text{جهت آزمون}&\text{مقدار }p\\\hline
\text{چپ‌طرفه}&G_0(t_{\rm obs})\\[2pt]
\text{راست‌طرفه}&1-G_0(t_{\rm obs})\\[2pt]
\text{دوطرفه با دو دُم برابر}&2\min\{G_0(t_{\rm obs}),1-G_0(t_{\rm obs})\}
\end{array}
\]
برای نرمال و استیودنتِ متقارن، مقدار دوطرفه را می‌توان به صورت $2[1-G_0(|t_{\rm obs}|)]$ نوشت. برای توزیع‌های نامتقارن، تعریف دوطرفه به شیوهٔ انتخاب ناحیهٔ رد وابسته است؛ فرمول جدول به آزمون دو دُم برابر مربوط می‌شود.
\begin{result}{قاعدهٔ تصمیم بر اساس مقدار $p$}
\[
p\le\alpha\quad\Longrightarrow\quad\text{رد }H_0,
\qquad
p>\alpha\quad\Longrightarrow\quad\text{رد نکردن }H_0.
\]
در مدل پیوسته، انتخاب $<$ یا $\le$ در مرز، احتمال خطا را تغییر نمی‌دهد؛ قرارداد تصمیم این جزوه $\le$ است.
\end{result}
\begin{example}{میانگین ضخامت محصول جدید}
ضخامت محصول قدیمی میانگین $500$ و انحراف معیار معلوم $20$ دارد. از محصول جدید، نمونه‌ای با $n=25$ و میانگین $508$ به دست آمده است. برای بررسی افزایش میانگین،
\[H_0:\mu=500,\qquad H_1:\mu>500,\qquad
z_{\rm obs}=\frac{508-500}{20/\sqrt{25}}=2.\]
با فرض مدل نرمال برای جامعه، یا مناسب بودن تقریب میانگین نمونه،
\[
p=\Prob(Z\ge2)=1-\Phi(2)=\boxed{0.0227501}.
\]
در سطح $\alpha=0.05$، چون $p<0.05$، فرض صفر رد می‌شود؛ در سطح $\alpha=0.01$، چون $p>0.01$، رد نمی‌شود.

همین نتیجه با مقادیر بحرانی به دست می‌آید: $z_{0.95}=1.64485<2$ ولی $z_{0.99}=2.32635>2$.
\end{example}
\begin{center}
\begin{LTR}
\begin{tikzpicture}[x=1.55cm,y=4cm]
\fill[teal!22] plot[domain=2:3.8,samples=65] (\x,{exp(-\x*\x/2)/sqrt(2*pi)})--(3.8,0)--(2,0)--cycle;
\draw[navy,thick] plot[domain=-3.4:3.8,samples=130] (\x,{exp(-\x*\x/2)/sqrt(2*pi)});
\draw[muted,-{Latex}] (-3.5,0)--(4,0) node[right] {$z$};
\draw[gold,dashed] (1.64485,0)--(1.64485,.29);
\draw[teal,thick] (2,0)--(2,.15);
\draw[muted,dashed] (2.32635,0)--(2.32635,.38);
\node[gold,anchor=south east] at (1.64485,.29) {$z_{0.95}=1.645$};
\node[muted,anchor=south west] at (2.32635,.38) {$z_{0.99}=2.326$};
\node[below,teal] at (2,0) {$z_{\rm obs}=2$};
\draw[-{Latex},teal] (3,.19) node[above] {$p=0.02275$}--(2.6,.012);
\end{tikzpicture}
\end{LTR}
\figcap{بازطراحی شکل‌های صفحه‌های ۴۳ تا ۴۵: مقدار مشاهده‌شده از حد بحرانی ۵ درصد عبور کرده، اما از حد بحرانی ۱ درصد عبور نکرده است. مساحت سبزآبی، مقدار $p$ را نشان می‌دهد.}
\end{center}
\section{رابطهٔ آزمون دوطرفه و فاصلهٔ اطمینان}
\begin{result}{دو بیان یک تصمیم}
در روش‌های متناظر، آزمون دوطرفهٔ سطح $\alpha$، فرض صفر $\theta=\theta_0$ را دقیقاً زمانی رد می‌کند که $\theta_0$ در فاصلهٔ اطمینان $1-\alpha$ متناظر نباشد. این هم‌ارزی به یکسان بودن مدل، آماره و روش ساخت دو طرف وابسته است.
\end{result}
برای مثال، در آزمون میانگین با واریانس معلوم،
\[
\left|\frac{\bar x-\mu_0}{\sigma/\sqrt n}\right|>z_{1-\alpha/2}
\quad\Longleftrightarrow\quad
\mu_0\notin\left[\bar x\pm z_{1-\alpha/2}\frac\sigma{\sqrt n}\right].
\]
\textbf{مسیر پیشنهادی حل مسائل آزمون.} ابتدا پارامتر و جهت اثر را مشخص کنید؛ فرض‌ها و سطح $\alpha$ را بنویسید؛ سپس شرایط روش و خطای استاندارد را تعیین کنید. آماره و مقدار $p$ یا حد بحرانی را محاسبه کنید و در پایان، تصمیم را در زبان مسئله همراه با جهت و اندازهٔ اثر گزارش دهید.

\appendix
\chapter{مرور سریع و راهنمای تطبیق}
\begingroup\small\linespread{1.0}\selectfont
\section{انتخاب روش برای میانگین‌ها}
\begin{center}\small
\renewcommand{\arraystretch}{1.15}
\begin{tabular}{>{\raggedleft\arraybackslash}p{57mm} >{\centering\arraybackslash}p{59mm} >{\centering\arraybackslash}p{24mm}}
\toprule
مسئله & خطای استاندارد & توزیع مرجع\\
\midrule
یک میانگین؛ واریانس معلوم & $\sigma/\sqrt n$ & $Z$\\
یک میانگین نرمال؛ واریانس نامعلوم & $S/\sqrt n$ & $t_{n-1}$\\
اختلاف دو میانگین؛ واریانس‌ها معلوم & $\sqrt{\sigma_1^2/n_1+\sigma_2^2/n_2}$ & $Z$\\
دو نمونهٔ مستقل؛ واریانس نامعلوم مشترک & $S_p\sqrt{1/n_1+1/n_2}$ & $t_{n_1+n_2-2}$\\
دو نمونهٔ مستقل؛ ولچ & $\sqrt{S_1^2/n_1+S_2^2/n_2}$ & $t_\nu$ تقریبی\\
داده‌های زوجی & $S_D/\sqrt n$ & $t_{n-1}$\\
\bottomrule
\end{tabular}
\end{center}
برای ساخت بازهٔ دوطرفه، «برآورد نقطه‌ای $\pm$ چندک $1-\alpha/2$ ضرب‌در خطای استاندارد» را به کار ببرید. برای آزمون، «برآورد منهای مقدار فرض صفر، تقسیم بر خطای استاندارد» را با چندک مناسب جهت آزمون مقایسه کنید.
\section{چند تمایز که از خطا جلوگیری می‌کند}
\begin{important}{پیش از جای‌گذاری عددها}
\textbf{پارامتر دوم نرمال}: واریانس است؛ برای به دست آوردن $\sigma$ ریشه بگیرید.

\textbf{سطح و جهت آزمون}: دوطرفهٔ ۵ درصد به چندک $0.975$ نیاز دارد؛ یک‌طرفهٔ ۵ درصد به چندک $0.95$.

\textbf{استقلال و زوجی بودن}: نمونهٔ مستقل و زوج‌های مرتبط، خطاهای استاندارد متفاوتی دارند.

\textbf{واریانس نمونه}: $S^2$ با مخرج $n-1$، درحالی‌که برآورد حول میانگین معلوم با مخرج $n$ است.

\textbf{نسبت‌ها}: در آزمون تک‌نسبتی، $p_0$؛ در فاصلهٔ والد، $\widehat p$؛ در آزمون برابری دو نسبت، نسبت تجمیعی را به کار ببرید.

\textbf{نااریبی و سازگاری}: اولی دربارهٔ امید برآوردگر در حجم نمونهٔ معین است؛ دومی دربارهٔ رفتار دنبالهٔ برآوردگرها با رشد حجم نمونه.
\end{important}
\clearpage
\section{مهم‌ترین بازسازی‌ها و اصلاحات علمی}
موارد زیر تغییرهای مؤثر بر معنا یا نتیجه‌اند؛ بقیهٔ تغییرها عمدتاً یکدست‌سازی نمادگذاری، نگارش و تکمیل مراحل محاسبه هستند.
\begin{enumerate}\small\setlength{\itemsep}{2pt}\setlength{\parsep}{0pt}
\item \textbf{صفحه‌های ۱ و ۲:} روابط چگالی به توزیع‌های مطلقاً پیوسته محدود شد و قیود وجود گشتاورها افزوده شد.
\item \textbf{صفحهٔ ۴:} رابطهٔ درست بازگشتی تابع گاما، $\Gamma(\alpha+1)=\alpha\Gamma(\alpha)$، نوشته شد و قرارداد نرخ مشخص شد.
\item \textbf{صفحه‌های ۶ تا ۸:} اثبات اریبی واریانس با مخرج $n$، حل ترکیب بهینه و شرایط کافی سازگاری تکمیل شد. مثال دوجمله‌ای با آمارهٔ نامعمول، همراه با خوانش استفاده‌شده حفظ شد.
\item \textbf{صفحه‌های ۱۰ و ۱۲ تا ۱۳:} استقلال لازم برای ترکیب نرمال و ساخت کای‌دو و استیودنت تصریح شد؛ استدلال درجات آزادی و درجهٔ صحیح $n-1$ اصلاح شد.
\item \textbf{صفحهٔ ۱۴:} حجم نمونهٔ مورد نیاز به $62$ گرد شد؛ حاشیهٔ خطا به‌عنوان کران احتمالی، و نه تضمین قطعی، تعبیر شد.
\item \textbf{صفحهٔ ۱۶:} برای واریانس دوم، خوانش $60$ انتخاب و نتیجهٔ خوانش جایگزین $90$ نیز گزارش شد.
\item \textbf{صفحهٔ ۱۸:} روابط ولچ به‌عنوان تقریب معرفی شد.
\item \textbf{صفحهٔ ۲۲:} بررسی برابری میانگین‌ها از بررسی برابری واریانس‌ها تفکیک شد.
\item \textbf{صفحه‌های ۲۳ و ۳۹:} فرض نرمال بودن اختلاف‌های زوجی به‌درستی بیان شد.
\item \textbf{صفحهٔ ۲۴:} واریانس اختلاف‌ها $8.7$ و فاصلهٔ اطمینان آن $[-2.862,4.462]$ محاسبه شد؛ حذف جملهٔ $n\bar d^2$ در دست‌نویس اصلاح شد.
\item \textbf{صفحه‌های ۲۵ و ۲۶:} فاصله‌های نسبت به‌عنوان روش تقریبی معرفی و محدودیت نمونهٔ دارای پنج موفقیت مشخص شد.
\item \textbf{صفحهٔ ۳۱:} در مثال با آمارهٔ $2$، حجم نمونهٔ $25$ از مراحل حل بازسازی شد.
\item \textbf{صفحهٔ ۳۲:} آزمون دوطرفهٔ نمرات ریاضی به چندک $t_{0.975,8}=2.306$ اصلاح شد؛ تصمیم صحیح در سطح ۵ درصد، رد نکردن فرض صفر است.
\item \textbf{صفحه‌های ۳۳ و ۳۴:} جهت تمام اختلاف میانگین‌ها به «اول منهای دوم» یکدست شد؛ جهت نامساوی‌ها نیز متناظر تنظیم شد.
\item \textbf{صفحه‌های ۴۳ تا ۴۵:} مقدار دقیق‌تر $p=0.0227501$ نوشته شد و تصمیم‌های سطح ۵ و ۱ درصد از هم تفکیک شدند. تفسیر «رد نکردن» جایگزین «اثبات یا قبول قطعی» شد.
\end{enumerate}
\clearpage
\section{نقشهٔ تطبیق با فایل ۴۵ صفحه‌ای}
شماره‌های ستون اول، شمارهٔ صفحات PDF ارسالی هستند؛ ممکن است با شماره‌های دست‌نویس گوشهٔ برگه‌ها یکسان نباشند. این جدول امکان بازبینی پوشش همهٔ صفحات را فراهم می‌کند.
\begin{center}\small
\renewcommand{\arraystretch}{1.2}
\begin{longtable}{>{\centering\arraybackslash}p{22mm} p{113mm}}
\toprule
صفحات مبنا & محل و موضوع بازنویسی\\
\midrule\endfirsthead
\toprule صفحات مبنا & محل و موضوع بازنویسی\\\midrule\endhead
۱–۲ & فصل اول: تعریف‌ها، تابع احتمال و چگالی، توزیع مشترک و امید ریاضی\\
۳–۴ & فصل اول: واریانس، کوواریانس، میانگین نمونه و توزیع‌های مهم\\
۵ & فصل دوم: جامعه، نمونه، آماره و مثال توزیع نمونه‌ای میانگین\\
۶ & فصل دوم: نااریبی، واریانس نمونه، خطای میانگین مربعات و کارایی\\
۷–۸ & فصل دوم: مثال‌های کارایی، ترکیب خطی و سازگاری؛ ارجاع‌های کتاب\\
۹–۱۱ & فصل سوم: فاصلهٔ اطمینان، کمیت محوری و میانگین با واریانس معلوم\\
۱۲–۱۳ & فصل سوم: کای‌دو و استیودنت؛ میانگین با واریانس نامعلوم\\
۱۴ & فصل سوم: حاشیهٔ خطا و تعیین حجم نمونه\\
۱۵ & فصل سوم: قضیهٔ حد مرکزی و اختلاف دو میانگین\\
۱۶–۱۸ & فصل سوم: اختلاف میانگین با واریانس‌های معلوم، مشترک و ولچ\\
۱۹–۲۰ & فصل سوم: فاصلهٔ واریانس با میانگین معلوم و نامعلوم\\
۲۰–۲۲ & فصل سوم: توزیع $F$، نسبت واریانس‌ها و خاصیت معکوس چندک\\
۲۳–۲۴ & فصل سوم: داده‌های زوجی و مثال نمرهٔ دو درس\\
۲۴–۲۵ & فصل سوم: فاصلهٔ اطمینان یک نسبت و مثال قطعات معیوب\\
۲۶ & فصل سوم: اختلاف دو نسبت و مثال نظر دانشجویان\\
۲۷–۲۸ & فصل چهارم: فرض‌ها، خطاها، توان و مثال سکه\\
۲۹–۳۱ & فصل چهارم: سه جهت آزمون میانگین با واریانس معلوم\\
۳۲ & فصل چهارم: آزمون استیودنت تک‌نمونه‌ای و مثال نمرات ریاضی\\
۳۳–۳۴ & فصل چهارم: آزمون اختلاف میانگین‌های مستقل\\
۳۵–۳۶ & فصل چهارم: آزمون واریانس یک جامعهٔ نرمال\\
۳۷–۳۸ & فصل چهارم: آزمون برابری واریانس دو جامعهٔ نرمال\\
۳۹ & فصل چهارم: آزمون میانگین در داده‌های زوجی\\
۴۰ & فصل چهارم: آزمون یک نسبت\\
۴۱–۴۲ & فصل چهارم: آزمون برابری دو نسبت و برآورد تجمیعی\\
۴۳–۴۵ & فصل چهارم: مقدار احتمالی، مثال محصول جدید و تصمیم در دو سطح\\
\bottomrule
\end{longtable}
\end{center}
\endgroup
\end{document}
